%% SPDX-License-Identifier: CC-BY-4.0
%% Copyright (C) 2026  Alin-Adrian Anton <alin.anton@upt.ro>
%%
%% data-diode-handbook.tex -- the reference document: threat model, software
%% diodes (bench only), hardware diodes, optical budgets, commissioning, and
% integration with the opc-diode tools.
%% Build:  make handbook   (or  latexmk -pdf data-diode-handbook.tex)

\documentclass[11pt,a4paper]{article}
\newcommand{\projectname}{opc-diode}
\newcommand{\doctitleshort}{The Data Diode Handbook}
%% SPDX-License-Identifier: CC-BY-4.0
%% Copyright (C) 2026  Alin-Adrian Anton <alin.anton@upt.ro>
%%
%% Shared with the sibling diode projects.  The canonical copy lives in
%% log-diode/docs/tutorial/ -- fix it there and re-copy, so opc-diode,
%% gnat-lt-pro, wal-postgresqldb and wal-mysqldb stay in step.
%% common.tex -- shared preamble for the diode tutorial documents.
%% Input from the preamble of both diode-quickstart.tex and
%% data-diode-handbook.tex.  Everything here must work with a stock
%% TeX Live / texlive-collection-latexrecommended install; no external
%% images, all figures are TikZ.

\usepackage[T1]{fontenc}
\usepackage{lmodern}
\usepackage{microtype}
\usepackage[a4paper,margin=25mm,headheight=15pt]{geometry}
\usepackage{amsmath,amssymb}
\usepackage{xcolor}
\usepackage{graphicx}      % \reflectbox (\copyleft), \resizebox for wide figures
\usepackage{booktabs,tabularx,longtable,array}
\usepackage{enumitem}
\usepackage{caption}
\usepackage{fancyhdr}
\usepackage{titlesec}
\usepackage{listings}
\usepackage{tikz}
\usetikzlibrary{positioning,arrows.meta,shapes.geometric,calc,fit,backgrounds}
\usepackage{tcolorbox}
\tcbuselibrary{skins,breakable}
\usepackage[hidelinks,pdfusetitle]{hyperref}

%% ---------------------------------------------------------------- colours
\definecolor{ldblue}  {HTML}{1F4E79}
\definecolor{ldred}   {HTML}{A61B1B}
\definecolor{ldgreen} {HTML}{1E6F3C}
\definecolor{ldamber} {HTML}{9A6700}
\definecolor{ldgray}  {HTML}{5A5A5A}
\definecolor{codebg}  {HTML}{F6F6F4}

%% ---------------------------------------------------------------- headers
\pagestyle{fancy}
\fancyhf{}
\fancyhead[L]{\footnotesize\textsc{\doctitleshort}}
\fancyhead[R]{\footnotesize\thepage}
\fancyfoot[C]{\footnotesize \projectname{} diode tutorial --- CC BY 4.0}
\renewcommand{\headrulewidth}{0.4pt}
\renewcommand{\footrulewidth}{0pt}

%% ---------------------------------------------------------------- listings
\lstset{
  basicstyle=\ttfamily\footnotesize,
  backgroundcolor=\color{codebg},
  frame=single,
  rulecolor=\color{ldgray!40},
  framesep=4pt,
  xleftmargin=6pt,
  xrightmargin=2pt,
  breaklines=true,
  breakatwhitespace=false,
  postbreak=\mbox{\textcolor{ldgray}{$\hookrightarrow$}\space},
  showstringspaces=false,
  keepspaces=true,
  columns=fullflexible,
  commentstyle=\color{ldgreen}\itshape,
  keywordstyle=\color{ldblue}\bfseries,
  stringstyle=\color{ldred},
  captionpos=b,
  aboveskip=8pt,
  belowskip=6pt,
  tabsize=2,
}

\lstdefinestyle{shell}{
  language=bash,
  morekeywords={ip,nft,iptables,ip6tables,pfctl,sysctl,systemctl,tcpdump,arp,
                ethtool,openssl,hping3,arping,nc,socat},
}
\lstdefinestyle{conf}{
  language=,
  morecomment=[l]{\#},
}
\lstdefinestyle{nft}{
  language=,
  morecomment=[l]{\#},
  morekeywords={table,chain,rule,type,hook,priority,policy,iif,oif,iifname,
                oifname,udp,dport,daddr,saddr,accept,drop,filter,ingress,
                input,output,forward,inet,netdev,ip,ip6,meta,counter,limit,log},
  keywordstyle=\color{ldblue}\bfseries,
}
\lstdefinestyle{pf}{
  language=,
  morecomment=[l]{\#},
  morekeywords={set,block,pass,in,out,quick,on,proto,from,to,port,all,
                no-state,state,antispoof,skip,block-policy,scrub,keep},
  keywordstyle=\color{ldblue}\bfseries,
}
\lstdefinestyle{ada}{
  language=Ada,
  morekeywords={SPARK_Mode,Pre,Post,Loop_Invariant},
}
\lstdefinestyle{c}{
  language=C,
  basicstyle=\ttfamily\scriptsize,
  morekeywords={ssize_t,size_t,sig_atomic_t,uint8_t,IFF_TAP,IFF_NO_PI,
                TUNSETIFF,POLLIN,O_RDWR,O_CLOEXEC,EINTR,EAGAIN,IFNAMSIZ},
  numbers=left,
  numberstyle=\ttfamily\tiny\color{ldgray!70},
  numbersep=7pt,
}

%% ------------------------------------------------------------- admonitions
\newtcolorbox{warnbox}[1][Warning]{%
  breakable, enhanced, colback=ldred!4, colframe=ldred,
  boxrule=0.9pt, left=6pt, right=6pt, top=5pt, bottom=5pt,
  fonttitle=\bfseries\sffamily\small, title={#1}}

\newtcolorbox{labbox}[1][Laboratory use only]{%
  breakable, enhanced, colback=ldamber!6, colframe=ldamber,
  boxrule=0.9pt, left=6pt, right=6pt, top=5pt, bottom=5pt,
  fonttitle=\bfseries\sffamily\small, title={#1}}

\newtcolorbox{notebox}[1][Note]{%
  breakable, enhanced, colback=ldblue!4, colframe=ldblue,
  boxrule=0.7pt, left=6pt, right=6pt, top=5pt, bottom=5pt,
  fonttitle=\bfseries\sffamily\small, title={#1}}

\newtcolorbox{tipbox}[1][In practice]{%
  breakable, enhanced, colback=ldgreen!4, colframe=ldgreen,
  boxrule=0.7pt, left=6pt, right=6pt, top=5pt, bottom=5pt,
  fonttitle=\bfseries\sffamily\small, title={#1}}

%% ------------------------------------------------------------------ macros
%% Copyleft (U+1F12F).  No Latin Modern / T1 glyph exists and neither
%% cclicenses nor ccicons is in a base TeX Live, so mirror the copyright
%% sign -- which is exactly what the symbol is.  \mathord keeps the spacing
%% sane when it lands next to text.
\newcommand{\copyleft}{\reflectbox{\copyright}}

\newcommand{\cmd}[1]{\texttt{#1}}
\newcommand{\file}[1]{\texttt{#1}}
\newcommand{\prog}[1]{\textsf{#1}}
%% Renewed per project, immediately after %% SPDX-License-Identifier: CC-BY-4.0
%% Copyright (C) 2026  Alin-Adrian Anton <alin.anton@upt.ro>
%%
%% Shared with the sibling diode projects.  The canonical copy lives in
%% log-diode/docs/tutorial/ -- fix it there and re-copy, so opc-diode,
%% gnat-lt-pro, wal-postgresqldb and wal-mysqldb stay in step.
%% common.tex -- shared preamble for the diode tutorial documents.
%% Input from the preamble of both diode-quickstart.tex and
%% data-diode-handbook.tex.  Everything here must work with a stock
%% TeX Live / texlive-collection-latexrecommended install; no external
%% images, all figures are TikZ.

\usepackage[T1]{fontenc}
\usepackage{lmodern}
\usepackage{microtype}
\usepackage[a4paper,margin=25mm,headheight=15pt]{geometry}
\usepackage{amsmath,amssymb}
\usepackage{xcolor}
\usepackage{graphicx}      % \reflectbox (\copyleft), \resizebox for wide figures
\usepackage{booktabs,tabularx,longtable,array}
\usepackage{enumitem}
\usepackage{caption}
\usepackage{fancyhdr}
\usepackage{titlesec}
\usepackage{listings}
\usepackage{tikz}
\usetikzlibrary{positioning,arrows.meta,shapes.geometric,calc,fit,backgrounds}
\usepackage{tcolorbox}
\tcbuselibrary{skins,breakable}
\usepackage[hidelinks,pdfusetitle]{hyperref}

%% ---------------------------------------------------------------- colours
\definecolor{ldblue}  {HTML}{1F4E79}
\definecolor{ldred}   {HTML}{A61B1B}
\definecolor{ldgreen} {HTML}{1E6F3C}
\definecolor{ldamber} {HTML}{9A6700}
\definecolor{ldgray}  {HTML}{5A5A5A}
\definecolor{codebg}  {HTML}{F6F6F4}

%% ---------------------------------------------------------------- headers
\pagestyle{fancy}
\fancyhf{}
\fancyhead[L]{\footnotesize\textsc{\doctitleshort}}
\fancyhead[R]{\footnotesize\thepage}
\fancyfoot[C]{\footnotesize \projectname{} diode tutorial --- CC BY 4.0}
\renewcommand{\headrulewidth}{0.4pt}
\renewcommand{\footrulewidth}{0pt}

%% ---------------------------------------------------------------- listings
\lstset{
  basicstyle=\ttfamily\footnotesize,
  backgroundcolor=\color{codebg},
  frame=single,
  rulecolor=\color{ldgray!40},
  framesep=4pt,
  xleftmargin=6pt,
  xrightmargin=2pt,
  breaklines=true,
  breakatwhitespace=false,
  postbreak=\mbox{\textcolor{ldgray}{$\hookrightarrow$}\space},
  showstringspaces=false,
  keepspaces=true,
  columns=fullflexible,
  commentstyle=\color{ldgreen}\itshape,
  keywordstyle=\color{ldblue}\bfseries,
  stringstyle=\color{ldred},
  captionpos=b,
  aboveskip=8pt,
  belowskip=6pt,
  tabsize=2,
}

\lstdefinestyle{shell}{
  language=bash,
  morekeywords={ip,nft,iptables,ip6tables,pfctl,sysctl,systemctl,tcpdump,arp,
                ethtool,openssl,hping3,arping,nc,socat},
}
\lstdefinestyle{conf}{
  language=,
  morecomment=[l]{\#},
}
\lstdefinestyle{nft}{
  language=,
  morecomment=[l]{\#},
  morekeywords={table,chain,rule,type,hook,priority,policy,iif,oif,iifname,
                oifname,udp,dport,daddr,saddr,accept,drop,filter,ingress,
                input,output,forward,inet,netdev,ip,ip6,meta,counter,limit,log},
  keywordstyle=\color{ldblue}\bfseries,
}
\lstdefinestyle{pf}{
  language=,
  morecomment=[l]{\#},
  morekeywords={set,block,pass,in,out,quick,on,proto,from,to,port,all,
                no-state,state,antispoof,skip,block-policy,scrub,keep},
  keywordstyle=\color{ldblue}\bfseries,
}
\lstdefinestyle{ada}{
  language=Ada,
  morekeywords={SPARK_Mode,Pre,Post,Loop_Invariant},
}
\lstdefinestyle{c}{
  language=C,
  basicstyle=\ttfamily\scriptsize,
  morekeywords={ssize_t,size_t,sig_atomic_t,uint8_t,IFF_TAP,IFF_NO_PI,
                TUNSETIFF,POLLIN,O_RDWR,O_CLOEXEC,EINTR,EAGAIN,IFNAMSIZ},
  numbers=left,
  numberstyle=\ttfamily\tiny\color{ldgray!70},
  numbersep=7pt,
}

%% ------------------------------------------------------------- admonitions
\newtcolorbox{warnbox}[1][Warning]{%
  breakable, enhanced, colback=ldred!4, colframe=ldred,
  boxrule=0.9pt, left=6pt, right=6pt, top=5pt, bottom=5pt,
  fonttitle=\bfseries\sffamily\small, title={#1}}

\newtcolorbox{labbox}[1][Laboratory use only]{%
  breakable, enhanced, colback=ldamber!6, colframe=ldamber,
  boxrule=0.9pt, left=6pt, right=6pt, top=5pt, bottom=5pt,
  fonttitle=\bfseries\sffamily\small, title={#1}}

\newtcolorbox{notebox}[1][Note]{%
  breakable, enhanced, colback=ldblue!4, colframe=ldblue,
  boxrule=0.7pt, left=6pt, right=6pt, top=5pt, bottom=5pt,
  fonttitle=\bfseries\sffamily\small, title={#1}}

\newtcolorbox{tipbox}[1][In practice]{%
  breakable, enhanced, colback=ldgreen!4, colframe=ldgreen,
  boxrule=0.7pt, left=6pt, right=6pt, top=5pt, bottom=5pt,
  fonttitle=\bfseries\sffamily\small, title={#1}}

%% ------------------------------------------------------------------ macros
%% Copyleft (U+1F12F).  No Latin Modern / T1 glyph exists and neither
%% cclicenses nor ccicons is in a base TeX Live, so mirror the copyright
%% sign -- which is exactly what the symbol is.  \mathord keeps the spacing
%% sane when it lands next to text.
\newcommand{\copyleft}{\reflectbox{\copyright}}

\newcommand{\cmd}[1]{\texttt{#1}}
\newcommand{\file}[1]{\texttt{#1}}
\newcommand{\prog}[1]{\textsf{#1}}
%% Renewed per project, immediately after %% SPDX-License-Identifier: CC-BY-4.0
%% Copyright (C) 2026  Alin-Adrian Anton <alin.anton@upt.ro>
%%
%% Shared with the sibling diode projects.  The canonical copy lives in
%% log-diode/docs/tutorial/ -- fix it there and re-copy, so opc-diode,
%% gnat-lt-pro, wal-postgresqldb and wal-mysqldb stay in step.
%% common.tex -- shared preamble for the diode tutorial documents.
%% Input from the preamble of both diode-quickstart.tex and
%% data-diode-handbook.tex.  Everything here must work with a stock
%% TeX Live / texlive-collection-latexrecommended install; no external
%% images, all figures are TikZ.

\usepackage[T1]{fontenc}
\usepackage{lmodern}
\usepackage{microtype}
\usepackage[a4paper,margin=25mm,headheight=15pt]{geometry}
\usepackage{amsmath,amssymb}
\usepackage{xcolor}
\usepackage{graphicx}      % \reflectbox (\copyleft), \resizebox for wide figures
\usepackage{booktabs,tabularx,longtable,array}
\usepackage{enumitem}
\usepackage{caption}
\usepackage{fancyhdr}
\usepackage{titlesec}
\usepackage{listings}
\usepackage{tikz}
\usetikzlibrary{positioning,arrows.meta,shapes.geometric,calc,fit,backgrounds}
\usepackage{tcolorbox}
\tcbuselibrary{skins,breakable}
\usepackage[hidelinks,pdfusetitle]{hyperref}

%% ---------------------------------------------------------------- colours
\definecolor{ldblue}  {HTML}{1F4E79}
\definecolor{ldred}   {HTML}{A61B1B}
\definecolor{ldgreen} {HTML}{1E6F3C}
\definecolor{ldamber} {HTML}{9A6700}
\definecolor{ldgray}  {HTML}{5A5A5A}
\definecolor{codebg}  {HTML}{F6F6F4}

%% ---------------------------------------------------------------- headers
\pagestyle{fancy}
\fancyhf{}
\fancyhead[L]{\footnotesize\textsc{\doctitleshort}}
\fancyhead[R]{\footnotesize\thepage}
\fancyfoot[C]{\footnotesize \projectname{} diode tutorial --- CC BY 4.0}
\renewcommand{\headrulewidth}{0.4pt}
\renewcommand{\footrulewidth}{0pt}

%% ---------------------------------------------------------------- listings
\lstset{
  basicstyle=\ttfamily\footnotesize,
  backgroundcolor=\color{codebg},
  frame=single,
  rulecolor=\color{ldgray!40},
  framesep=4pt,
  xleftmargin=6pt,
  xrightmargin=2pt,
  breaklines=true,
  breakatwhitespace=false,
  postbreak=\mbox{\textcolor{ldgray}{$\hookrightarrow$}\space},
  showstringspaces=false,
  keepspaces=true,
  columns=fullflexible,
  commentstyle=\color{ldgreen}\itshape,
  keywordstyle=\color{ldblue}\bfseries,
  stringstyle=\color{ldred},
  captionpos=b,
  aboveskip=8pt,
  belowskip=6pt,
  tabsize=2,
}

\lstdefinestyle{shell}{
  language=bash,
  morekeywords={ip,nft,iptables,ip6tables,pfctl,sysctl,systemctl,tcpdump,arp,
                ethtool,openssl,hping3,arping,nc,socat},
}
\lstdefinestyle{conf}{
  language=,
  morecomment=[l]{\#},
}
\lstdefinestyle{nft}{
  language=,
  morecomment=[l]{\#},
  morekeywords={table,chain,rule,type,hook,priority,policy,iif,oif,iifname,
                oifname,udp,dport,daddr,saddr,accept,drop,filter,ingress,
                input,output,forward,inet,netdev,ip,ip6,meta,counter,limit,log},
  keywordstyle=\color{ldblue}\bfseries,
}
\lstdefinestyle{pf}{
  language=,
  morecomment=[l]{\#},
  morekeywords={set,block,pass,in,out,quick,on,proto,from,to,port,all,
                no-state,state,antispoof,skip,block-policy,scrub,keep},
  keywordstyle=\color{ldblue}\bfseries,
}
\lstdefinestyle{ada}{
  language=Ada,
  morekeywords={SPARK_Mode,Pre,Post,Loop_Invariant},
}
\lstdefinestyle{c}{
  language=C,
  basicstyle=\ttfamily\scriptsize,
  morekeywords={ssize_t,size_t,sig_atomic_t,uint8_t,IFF_TAP,IFF_NO_PI,
                TUNSETIFF,POLLIN,O_RDWR,O_CLOEXEC,EINTR,EAGAIN,IFNAMSIZ},
  numbers=left,
  numberstyle=\ttfamily\tiny\color{ldgray!70},
  numbersep=7pt,
}

%% ------------------------------------------------------------- admonitions
\newtcolorbox{warnbox}[1][Warning]{%
  breakable, enhanced, colback=ldred!4, colframe=ldred,
  boxrule=0.9pt, left=6pt, right=6pt, top=5pt, bottom=5pt,
  fonttitle=\bfseries\sffamily\small, title={#1}}

\newtcolorbox{labbox}[1][Laboratory use only]{%
  breakable, enhanced, colback=ldamber!6, colframe=ldamber,
  boxrule=0.9pt, left=6pt, right=6pt, top=5pt, bottom=5pt,
  fonttitle=\bfseries\sffamily\small, title={#1}}

\newtcolorbox{notebox}[1][Note]{%
  breakable, enhanced, colback=ldblue!4, colframe=ldblue,
  boxrule=0.7pt, left=6pt, right=6pt, top=5pt, bottom=5pt,
  fonttitle=\bfseries\sffamily\small, title={#1}}

\newtcolorbox{tipbox}[1][In practice]{%
  breakable, enhanced, colback=ldgreen!4, colframe=ldgreen,
  boxrule=0.7pt, left=6pt, right=6pt, top=5pt, bottom=5pt,
  fonttitle=\bfseries\sffamily\small, title={#1}}

%% ------------------------------------------------------------------ macros
%% Copyleft (U+1F12F).  No Latin Modern / T1 glyph exists and neither
%% cclicenses nor ccicons is in a base TeX Live, so mirror the copyright
%% sign -- which is exactly what the symbol is.  \mathord keeps the spacing
%% sane when it lands next to text.
\newcommand{\copyleft}{\reflectbox{\copyright}}

\newcommand{\cmd}[1]{\texttt{#1}}
\newcommand{\file}[1]{\texttt{#1}}
\newcommand{\prog}[1]{\textsf{#1}}
%% Renewed per project, immediately after \input{common}.
\newcommand{\amp}{\prog{sender}}
\newcommand{\deamp}{\prog{receiver}}
%% The two ends of the link, named by ROLE rather than by security level.
%% A diode can be deployed in either orientation -- ingesting into a
%% protected domain, or exporting out of one -- so the mechanics here are
%% written in terms of which end sends and which receives. Each project's
%% "Orientation" section maps these onto its own security levels.
\newcommand{\srcside}{\textsc{source}}
\newcommand{\dstside}{\textsc{destination}}
\newcommand{\nm}[1]{#1\,nm}
\newcommand{\dbm}[1]{#1\,dBm}
\newcommand{\dbb}[1]{#1\,dB}

%% ------------------------------------------------------------ tikz styles
\tikzset{
  host/.style     ={draw, rounded corners=2pt, minimum width=27mm,
                    minimum height=12mm, align=center, font=\small,
                    fill=ldblue!8, draw=ldblue!70},
  mc/.style       ={draw, minimum width=25mm, minimum height=13mm,
                    align=center, font=\scriptsize, fill=gray!8,
                    draw=ldgray},
  netbox/.style   ={draw, dashed, rounded corners=3pt, draw=ldgray!70,
                    inner sep=5mm},
  splitter/.style ={draw, diamond, aspect=1.5, inner sep=1.5pt,
                    align=center, font=\scriptsize, fill=ldgreen!12,
                    draw=ldgreen},
  fiber/.style    ={-{Latex[length=2.4mm,width=1.8mm]}, thick,
                    draw=ldgreen!65!black},
  copper/.style   ={-, thick, draw=ldblue!70!black},
  dataflow/.style ={-{Latex[length=3.4mm,width=2.6mm]}, line width=1.5pt,
                    draw=ldblue},
  blocked/.style  ={-{Latex[length=2.4mm]}, thick, draw=ldred, dashed},
  lbl/.style      ={font=\scriptsize\sffamily, inner sep=1.5pt},
  lblg/.style     ={font=\scriptsize\sffamily, inner sep=1.5pt,
                    text=ldgreen!55!black},
  lblr/.style     ={font=\scriptsize\sffamily, inner sep=1.5pt, text=ldred},
}

%% A red "no return path" marker, placed on a path with node[midway]{\xmark}.
\newcommand{\xmark}{\textcolor{ldred}{\sffamily\bfseries\large $\times$}}

%% ------------------------------------------------------------ title block
\titleformat{\section}{\normalfont\Large\bfseries\color{ldblue}}{\thesection}{0.7em}{}
\titleformat{\subsection}{\normalfont\large\bfseries\color{ldblue!85!black}}{\thesubsection}{0.6em}{}
\titleformat{\subsubsection}{\normalfont\normalsize\bfseries}{\thesubsubsection}{0.6em}{}
\setlist[itemize]{itemsep=2pt, topsep=3pt, parsep=0pt}
\setlist[enumerate]{itemsep=2pt, topsep=3pt, parsep=0pt}
\captionsetup{font=small, labelfont={bf,color=ldblue}}
.
\newcommand{\amp}{\prog{sender}}
\newcommand{\deamp}{\prog{receiver}}
%% The two ends of the link, named by ROLE rather than by security level.
%% A diode can be deployed in either orientation -- ingesting into a
%% protected domain, or exporting out of one -- so the mechanics here are
%% written in terms of which end sends and which receives. Each project's
%% "Orientation" section maps these onto its own security levels.
\newcommand{\srcside}{\textsc{source}}
\newcommand{\dstside}{\textsc{destination}}
\newcommand{\nm}[1]{#1\,nm}
\newcommand{\dbm}[1]{#1\,dBm}
\newcommand{\dbb}[1]{#1\,dB}

%% ------------------------------------------------------------ tikz styles
\tikzset{
  host/.style     ={draw, rounded corners=2pt, minimum width=27mm,
                    minimum height=12mm, align=center, font=\small,
                    fill=ldblue!8, draw=ldblue!70},
  mc/.style       ={draw, minimum width=25mm, minimum height=13mm,
                    align=center, font=\scriptsize, fill=gray!8,
                    draw=ldgray},
  netbox/.style   ={draw, dashed, rounded corners=3pt, draw=ldgray!70,
                    inner sep=5mm},
  splitter/.style ={draw, diamond, aspect=1.5, inner sep=1.5pt,
                    align=center, font=\scriptsize, fill=ldgreen!12,
                    draw=ldgreen},
  fiber/.style    ={-{Latex[length=2.4mm,width=1.8mm]}, thick,
                    draw=ldgreen!65!black},
  copper/.style   ={-, thick, draw=ldblue!70!black},
  dataflow/.style ={-{Latex[length=3.4mm,width=2.6mm]}, line width=1.5pt,
                    draw=ldblue},
  blocked/.style  ={-{Latex[length=2.4mm]}, thick, draw=ldred, dashed},
  lbl/.style      ={font=\scriptsize\sffamily, inner sep=1.5pt},
  lblg/.style     ={font=\scriptsize\sffamily, inner sep=1.5pt,
                    text=ldgreen!55!black},
  lblr/.style     ={font=\scriptsize\sffamily, inner sep=1.5pt, text=ldred},
}

%% A red "no return path" marker, placed on a path with node[midway]{\xmark}.
\newcommand{\xmark}{\textcolor{ldred}{\sffamily\bfseries\large $\times$}}

%% ------------------------------------------------------------ title block
\titleformat{\section}{\normalfont\Large\bfseries\color{ldblue}}{\thesection}{0.7em}{}
\titleformat{\subsection}{\normalfont\large\bfseries\color{ldblue!85!black}}{\thesubsection}{0.6em}{}
\titleformat{\subsubsection}{\normalfont\normalsize\bfseries}{\thesubsubsection}{0.6em}{}
\setlist[itemize]{itemsep=2pt, topsep=3pt, parsep=0pt}
\setlist[enumerate]{itemsep=2pt, topsep=3pt, parsep=0pt}
\captionsetup{font=small, labelfont={bf,color=ldblue}}
.
\newcommand{\amp}{\prog{sender}}
\newcommand{\deamp}{\prog{receiver}}
%% The two ends of the link, named by ROLE rather than by security level.
%% A diode can be deployed in either orientation -- ingesting into a
%% protected domain, or exporting out of one -- so the mechanics here are
%% written in terms of which end sends and which receives. Each project's
%% "Orientation" section maps these onto its own security levels.
\newcommand{\srcside}{\textsc{source}}
\newcommand{\dstside}{\textsc{destination}}
\newcommand{\nm}[1]{#1\,nm}
\newcommand{\dbm}[1]{#1\,dBm}
\newcommand{\dbb}[1]{#1\,dB}

%% ------------------------------------------------------------ tikz styles
\tikzset{
  host/.style     ={draw, rounded corners=2pt, minimum width=27mm,
                    minimum height=12mm, align=center, font=\small,
                    fill=ldblue!8, draw=ldblue!70},
  mc/.style       ={draw, minimum width=25mm, minimum height=13mm,
                    align=center, font=\scriptsize, fill=gray!8,
                    draw=ldgray},
  netbox/.style   ={draw, dashed, rounded corners=3pt, draw=ldgray!70,
                    inner sep=5mm},
  splitter/.style ={draw, diamond, aspect=1.5, inner sep=1.5pt,
                    align=center, font=\scriptsize, fill=ldgreen!12,
                    draw=ldgreen},
  fiber/.style    ={-{Latex[length=2.4mm,width=1.8mm]}, thick,
                    draw=ldgreen!65!black},
  copper/.style   ={-, thick, draw=ldblue!70!black},
  dataflow/.style ={-{Latex[length=3.4mm,width=2.6mm]}, line width=1.5pt,
                    draw=ldblue},
  blocked/.style  ={-{Latex[length=2.4mm]}, thick, draw=ldred, dashed},
  lbl/.style      ={font=\scriptsize\sffamily, inner sep=1.5pt},
  lblg/.style     ={font=\scriptsize\sffamily, inner sep=1.5pt,
                    text=ldgreen!55!black},
  lblr/.style     ={font=\scriptsize\sffamily, inner sep=1.5pt, text=ldred},
}

%% A red "no return path" marker, placed on a path with node[midway]{\xmark}.
\newcommand{\xmark}{\textcolor{ldred}{\sffamily\bfseries\large $\times$}}

%% ------------------------------------------------------------ title block
\titleformat{\section}{\normalfont\Large\bfseries\color{ldblue}}{\thesection}{0.7em}{}
\titleformat{\subsection}{\normalfont\large\bfseries\color{ldblue!85!black}}{\thesubsection}{0.6em}{}
\titleformat{\subsubsection}{\normalfont\normalsize\bfseries}{\thesubsubsection}{0.6em}{}
\setlist[itemize]{itemsep=2pt, topsep=3pt, parsep=0pt}
\setlist[enumerate]{itemsep=2pt, topsep=3pt, parsep=0pt}
\captionsetup{font=small, labelfont={bf,color=ldblue}}

\renewcommand{\amp}{\prog{od\_sender}}
\renewcommand{\deamp}{\prog{od\_receiver}}

\title{\vspace{-12mm}\textbf{The Data Diode Handbook}\\[2mm]
       \large Software and hardware unidirectional links,\\
       \large and the proven OPC UA relay that runs across them}
\author{Alin-Adrian Anton \\
        \small \texttt{alin.anton@upt.ro} \\
        \small Computer and Information Technology Department\\
        \small Politehnica University of Timi\c{s}oara}
\date{\small For the opc-diode project\\
      \small Diode build material shared across the sibling
      projects; canonical copy in log-diode}

\begin{document}
\maketitle
\thispagestyle{fancy}

\begin{abstract}
\noindent
A data diode enforces one-way information flow. This handbook covers how to
build one, how to prove you have built one, and how to run useful traffic
across it. It treats three families: \textbf{firewall-based diodes}
(\prog{iptables}, \prog{nftables}, \prog{pf}); a \textbf{unidirectional
layer-2 tunnel} built by bridging a TAP interface to each of two real
Ethernet segments and forwarding frames with a program that has no reverse
code path; and \textbf{physical diodes} --- commodity media converters between
networks, or a three-wire SPI link between two microcontrollers. The first two
are excellent bench instruments and are \emph{not} security boundaries; only
the third survives a compromise of the hosts it separates. For the optical
case it works through both published topologies --- two converters with a
fiber splitter on the transmitter's return path, and three converters in a
triangle --- with the power budgets, the attenuation they imply, and the
procurement reality that actually decides between them. It closes with the
transport problem a diode creates, and with \amp{}/\deamp{}: Ada/SPARK
daemons whose parsing and reconstruction path is machine-checked free of
run-time errors, because on a link with no feedback there is no second
chance.
\end{abstract}

\vspace{2mm}
{\small\tableofcontents}
\clearpage

%% =====================================================================
\section{Scope}

\subsection{What this document is}

The \emph{Sensors} paper \cite{anton2024} gives the essential recipe for
the diode itself in its \S2.2: static ARP, addressing,
and two media-converter topologies. This handbook expands that recipe into
something you can hand to an engineer, and adds three things the paper does
not cover:

\begin{enumerate}
\item \textbf{Software diodes.} The paper does not discuss them. Both the
      firewall rulesets of \S\ref{sec:soft} and the bridged layer-2 tunnel of
      \S\ref{sec:tap} are genuinely useful --- for development, CI and
      teaching --- and genuinely dangerous if mistaken for the real thing, so
      they get an honest chapter with an honest warning.
\item \textbf{The optical practicalities.} \S\ref{sec:budget} gives the power
      budgets, says when to reach for an attenuator, and is honest about the
      part that actually gates a splitter build --- an \nm{850} multimode
      splitter is a special order with a long lead time, while \nm{1310} ones
      are catalogue stock.
\item \textbf{Embedded scale.} Between two microcontrollers a media converter
      is absurd. \S\ref{sec:spi} builds the same guarantee from three wires of
      SPI, and explains why the sender --- not the receiver --- has to be the
      bus master.
\item \textbf{Commissioning.} How to \emph{prove} a link is one-way, rather
      than believe it.
\end{enumerate}

\subsection{How to read it}

The \emph{Quick Start} companion is the build sheet; this is the reference.
If you want a working diode today, read that first and come back here for the
reasoning. Sections \S\ref{sec:threat} and \S\ref{sec:commission} are the ones
to read before anything goes into production.

\subsection{Conventions}

Throughout, \srcside{} is the network data leaves and \dstside{} is the
network it arrives in. Information flows
\srcside{}~$\rightarrow$~\dstside{} only.

These are \emph{roles, not security levels}, and the distinction matters. A
diode is deployed in one of two orientations, and which one you are building
decides what the device is protecting:

\begin{itemize}
\item \textbf{Ingest} --- the source is the lower-security network and the
      destination the higher-security one. The diode protects the
      destination's \emph{confidentiality}: the untrusted side can send in but
      can never read back. This is the Bell--LaPadula arrangement
      \cite{bell1976}, and it is how log collection works --- a SOC ingests
      logs from a network where an attacker may already be.
\item \textbf{Export} --- the source is the higher-security network and the
      destination the lower-security one. The diode protects the source's
      \emph{integrity and reachability}: the protected side publishes outward,
      and nothing on the receiving network can reach in, write to it, or probe
      it. This is how a plant exports process data or a database ships its
      write-ahead log to a mirror.
\end{itemize}

Everything mechanical in this handbook --- the firewall rules, the media
converters, the SPI link, the commissioning tests --- is identical in both
orientations, which is why it is written in terms of source and destination.
\S\ref{sec:orient} states which orientation this project uses.

Addresses are \cmd{192.168.10.1} (\srcside{}) and \cmd{192.168.10.2}
(\dstside{}); the diode-facing interface is \cmd{enp1s0} on Linux and
\cmd{em0} on BSD.

%% =====================================================================
\section{Threat model and assurance tiers}
\label{sec:threat}

\subsection{The property being claimed}

A data diode makes one claim: \emph{no information can travel from
\dstside{} back to \srcside{}}. Everything else --- confidentiality,
integrity, reliability --- is built on top of that claim by the software
running across the link, not by the link itself.

That single claim is what both orientations rest on. In an ingest deployment
it means the untrusted sender can never read the protected receiver, which is
Bell--LaPadula confidentiality \cite{bell1976}. In an export deployment it
means the receiving network can never reach back into the protected sender,
which is an integrity and isolation property rather than a confidentiality
one. The hardware is the same; only what you are protecting changes.

\subsection{Orientation for this project}
\label{sec:orient}

\textbf{opc-diode is an export deployment.} The source is the
\emph{higher}-security network: a plant publishes OPC UA PubSub (UADP over
UDP) natively, and the diode carries those NetworkMessages out to subscribers
on a lower-security monitoring network. So \srcside{}~=~high security and
\dstside{}~=~low security here --- the reverse of an ingest deployment.

The property being bought is therefore not confidentiality of the receiver but
\textbf{integrity and reachability of the sender}: the monitoring network can
read process data, and can never reach back into the plant to write a
setpoint, probe a PLC, or pivot. PubSub helps here, being connectionless and
one-way by construction, so there is no session for the diode to break.

\begin{notebox}[What a diode does not do]
A diode says nothing about the \emph{truthfulness} of what crosses it. A
compromised device on \srcside{} can send whatever it likes, including forged
or replayed log lines. The diode guarantees only that the attacker cannot
\emph{see} or \emph{reach} \dstside{}. Authenticity of content is the
cryptography's job (\S\ref{sec:integrate}), and the authenticity of the
\emph{sender} depends on the security of the \srcside{} devices themselves.
\end{notebox}

\subsection{Three tiers, and where each fails}

\begin{table}[h]
\centering\small
\begin{tabularx}{\textwidth}{@{}p{28mm}p{34mm}X@{}}
\toprule
\textbf{Tier} & \textbf{Enforced by} & \textbf{How it fails} \\
\midrule
Application &
  A daemon that chooses not to reply &
  A bug, a configuration change, or code execution in the daemon. Offers
  essentially no assurance; not treated further here. \\[3pt]
Firewall &
  A ruleset in the kernel &
  Root compromise, \cmd{nft flush ruleset}, a configuration-management run, a
  kernel or conntrack bug, an interface rename. Recovers a full return path
  instantly and silently. \\[3pt]
Layer-2 tunnel (\S\ref{sec:tap}) &
  A forwarding program with no reverse code path &
  Replacing the binary, or bridging both NICs together. Stronger than a
  ruleset --- the reverse path is absent rather than forbidden --- but still
  one command away for root on the forwarding host. \\[3pt]
Physical\newline (optics, or SPI with MISO omitted) &
  There is no transmitter, and no conductor, pointed the wrong way &
  Someone rewires it, or fits the missing wire. Nothing software can do
  restores the path. \\
\bottomrule
\end{tabularx}
\caption{Assurance tiers. Only the optical tier survives a compromise of the
         hosts it separates.}
\end{table}

The distinction matters because of \emph{who} you are defending against. If
the adversary is a misconfigured application, a firewall diode helps. If the
adversary has code execution on the \srcside{} host --- which is the entire
premise of putting a diode between the log producers and the SOC --- then the
firewall is inside the attacker's blast radius and the guarantee evaporates.
As the paper puts it, the highest security diode enforces unidirectionality
from physical properties and is immune from software exploits
\cite{anton2024}.

\subsection{Residual risks even with optics}

\begin{itemize}
\item \textbf{Rewiring.} A person with physical access can restore a return
      path in seconds. This is the paper's ``unintentional insider'': an
      administrator who tidies the cables. Mitigation is procedural ---
      labelling, tamper-evident seals, locked enclosures, and the audit in
      \S\ref{sec:commission}.
\item \textbf{Covert channels.} A colluding insider can exfiltrate through
      channels the diode does not model: NIC LEDs \cite{etherled}, Ethernet
      cable emissions \cite{lantenna}, power lines \cite{powerhammer}, and
      memory-bus emissions picked up by a nearby phone \cite{gsmem}. A diode
      is not a TEMPEST control.
\item \textbf{Loss.} A one-way link cannot retransmit. Loss beyond the error
      code's capacity is unrecoverable. This is the honest failure mode of a
      diode and is a design constraint, not a defect.
\item \textbf{No diagnostics.} Nothing can report back. Failure is silent by
      construction; see \S\ref{sec:maint}.
\item \textbf{Availability.} Every converter is a point of failure, and a
      three-converter diode has three of them.
\end{itemize}

%% =====================================================================
\section{Why the transport looks the way it does}

Remove the return path and a surprising amount of ordinary networking stops
working. It is worth being explicit, because every one of these becomes a
configuration task.

\begin{table}[h]
\centering\small
\begin{tabularx}{\textwidth}{@{}p{30mm}X@{}}
\toprule
\textbf{What breaks} & \textbf{Consequence} \\
\midrule
TCP & No handshake, so no TCP at all. Everything must be UDP or a raw
       protocol. \\
ARP & The sender never learns the receiver's MAC. Must be configured
      statically --- this is why \file{/etc/ethers} appears in every recipe
      here. \\
ICMP & No \cmd{ping}, no path MTU discovery. Set the MTU by hand and keep
       datagrams comfortably under it. \\
DNS, NTP, DHCP & All require a reply. Use static addressing, static
       \file{/etc/hosts}, and an \dstside{}-local time source. \\
Retransmission & Nothing can be requested again. Redundancy has to be sent
       \emph{in advance}. \\
Key exchange & No Diffie--Hellman, no TLS. Keys are pre-shared out of band. \\
\bottomrule
\end{tabularx}
\end{table}

The last two shape the software. Since a lost datagram is lost forever, the
sender must decide up front how much redundancy to spend. The paper's answer
was a repeat factor: send each message $R$ times. The tools described in
\S\ref{sec:integrate} improve on it in two ways --- a Reed--Solomon erasure
code, so that any $K$ of $K+M$ fragments rebuild the message rather than
needing one intact copy; and interleaving, so that a burst of loss cannot
consume every copy of the same message.

%% =====================================================================
\section{Software diodes}
\label{sec:soft}

\begin{labbox}[This entire section is for the bench]
A firewall diode enforces direction with a mutable ruleset inside a kernel
that the attacker may be able to reach. \cmd{nft flush ruleset} is one
command. A configuration-management run that re-applies a site-wide policy is
one mistake. Treat everything below as a \textbf{laboratory and CI
instrument}: it lets you develop against \amp{}/\deamp{}, write integration
tests, and teach the topology without buying hardware. It is not a boundary
between two security domains. For that, build \S\ref{sec:hard}.
\end{labbox}

\subsection{The one principle}

Every ruleset in this section obeys a single negative rule:

\begin{center}
\emph{Nothing on the diode interface may create or consult connection state.}
\end{center}

A conntrack entry, an \cmd{ESTABLISHED,RELATED} accept, or a pf
\cmd{keep state} exists precisely to let the reply through. So does a TCP RST
and an ICMP port-unreachable --- both are packets travelling the wrong way,
and both leak information (they tell the \dstside{} side that something is
listening). Stateless, silent dropping is the only correct behaviour.

\begin{figure}[h]
\centering
\begin{tikzpicture}
  \node[host] (low) {SOURCE host\\\scriptsize 192.168.10.1};
  \node[host, right=45mm of low] (high) {DESTINATION host\\\scriptsize 192.168.10.2};
  \draw[dataflow] (low) -- node[lbl,above]{UDP 6514 only}
                           node[lbl,below]{no state, no conntrack} (high);
  \draw[blocked] (high.south) to[out=-150,in=-30]
        node[midway,fill=white,inner sep=1pt]{\xmark}
        node[lbl,below=3.5mm]{OUTPUT policy drop \emph{and} INPUT policy drop} (low.south);
  \node[lbl, above=9mm of $(low)!0.5!(high)$, align=center, text=ldamber!70!black]
       {\textbf{enforced by a ruleset --- revocable by whoever owns root}};
\end{tikzpicture}
\caption{The software diode. The asymmetry is real but it lives in mutable
         kernel state.}
\end{figure}

\subsection{Link-layer preparation (all platforms)}
\label{sec:linklayer}

This layer is needed for hardware diodes too --- with no return path, address
resolution has to be done by hand.

\subsubsection{Linux}

Record the peer's MAC in \file{/etc/ethers} on each host (\srcside{} records
\dstside{}'s, and vice versa):

\begin{lstlisting}[style=conf]
# /etc/ethers on the SOURCE host
bb:bb:bb:bb:bb:bb 192.168.10.2
\end{lstlisting}

\begin{lstlisting}[style=shell]
ip addr add 192.168.10.1/24 dev enp1s0     # .2/24 on the DESTINATION host
ip link set enp1s0 up
arp -f /etc/ethers                          # net-tools package
ip neigh show dev enp1s0                    # must report PERMANENT
\end{lstlisting}

The paper loads this from \file{/etc/rc.local} via a \cmd{rc-local.service}
shim \cite{anton2024}; a dedicated unit is cleaner on a modern system:

\begin{lstlisting}[style=conf,caption={\file{/etc/systemd/system/diode-link.service}}]
[Unit]
Description=Static ARP and addressing for the diode link
After=network-pre.target
Wants=network-pre.target

[Service]
Type=oneshot
RemainAfterExit=yes
ExecStart=/usr/sbin/arp -f /etc/ethers

[Install]
WantedBy=multi-user.target
\end{lstlisting}

Then quieten the interface:

\begin{lstlisting}[style=conf,caption={\file{/etc/sysctl.d/90-diode.conf}}]
# Never answer an ARP request on the diode NIC; never volunteer our address.
net.ipv4.conf.enp1s0.arp_ignore   = 8
net.ipv4.conf.enp1s0.arp_announce = 2

# IPv6 is chatty and inherently bidirectional (NDP, RA). Turn it off here.
net.ipv6.conf.enp1s0.disable_ipv6 = 1

net.ipv4.conf.enp1s0.accept_redirects    = 0
net.ipv4.conf.enp1s0.send_redirects      = 0
net.ipv4.conf.enp1s0.accept_source_route = 0
\end{lstlisting}

\begin{warnbox}[Reverse-path filtering]
Leave \cmd{rp\_filter} alone unless you know your topology needs it. In the
symmetric-subnet layout used here, strict reverse-path filtering
(\cmd{rp\_filter=1}) passes, because the route to the peer's source address
does point out of \cmd{enp1s0}. If you place the two sides in \emph{different}
subnets --- a common choice, since routing back is impossible anyway --- strict
RPF will silently discard every arriving datagram on the \dstside{} host.
That failure looks exactly like a broken diode. Set
\cmd{net.ipv4.conf.enp1s0.rp\_filter = 0} in that case and note it in the
commissioning record.
\end{warnbox}

If your distribution runs \prog{NetworkManager}, exclude the diode interface
so it does not re-apply DHCP or IPv6 autoconfiguration behind your back:

\begin{lstlisting}[style=conf,caption={\file{/etc/NetworkManager/conf.d/90-diode.conf}}]
[keyfile]
unmanaged-devices=interface-name:enp1s0
\end{lstlisting}

\subsubsection{BSD}

\begin{lstlisting}[style=shell,caption={OpenBSD}]
ifconfig em0 inet 192.168.10.1/24 up
arp -s 192.168.10.2 bb:bb:bb:bb:bb:bb permanent
# persist: /etc/hostname.em0 for the address, /etc/rc.local for the arp entry
\end{lstlisting}

\begin{lstlisting}[style=conf,caption={FreeBSD \file{/etc/rc.conf}}]
ifconfig_em0="inet 192.168.10.1 netmask 255.255.255.0"
static_arp_pairs="diodepeer"
static_arp_diodepeer="192.168.10.2 bb:bb:bb:bb:bb:bb"
pf_enable="YES"
pf_rules="/etc/pf.conf"
\end{lstlisting}

\subsection{nftables}

\prog{nftables} is the recommended Linux option, for one reason that
\prog{iptables} cannot match: the \cmd{netdev} family's \cmd{ingress} hook
sees \emph{every} frame, including ARP, LLDP and IPv6, before the IP stack
does. A single \cmd{policy drop} there makes the interface deaf at the link
layer.

\begin{lstlisting}[style=nft,caption={\file{/etc/nftables.conf} --- SOURCE host (transmit only)}]
#!/usr/sbin/nft -f
flush ruleset

define PEER = 192.168.10.2
define PORT = 6514

# Link-layer deafness. Note: the device name cannot be a variable in
# older nft releases, so it is spelled out here.
table netdev diodeguard {
    chain ingress {
        type filter hook ingress device enp1s0 priority -500; policy drop;
    }
}

table inet diode {
    chain input {
        type filter hook input priority filter; policy drop;
        iif lo accept
        iifname "enp1s0" counter drop
    }
    chain output {
        type filter hook output priority filter; policy drop;
        oif lo accept
        oifname "enp1s0" ip daddr $PEER udp dport $PORT counter accept
        oifname "enp1s0" counter drop
    }
    chain forward {
        type filter hook forward priority filter; policy drop;
    }
}
\end{lstlisting}

\begin{lstlisting}[style=nft,caption={\file{/etc/nftables.conf} --- DESTINATION host (receive only)}]
#!/usr/sbin/nft -f
flush ruleset

define PEER = 192.168.10.1
define PORT = 6514

table inet diode {
    chain input {
        type filter hook input priority filter; policy drop;
        iif lo accept
        iifname "enp1s0" ip saddr $PEER udp dport $PORT counter accept
        iifname "enp1s0" counter drop
    }
    chain output {
        type filter hook output priority filter; policy drop;
        oif lo accept
        # Absolutely nothing leaves by the diode NIC.
        oifname "enp1s0" counter drop
    }
    chain forward {
        type filter hook forward priority filter; policy drop;
    }
}
\end{lstlisting}

The \cmd{counter} keywords are not decoration. After a test run,
\cmd{nft list ruleset} shows how many packets each rule matched; a non-zero
counter on the \dstside{} host's output-drop rule is direct evidence that
something tried to talk backwards.

\begin{lstlisting}[style=shell]
systemctl enable --now nftables
nft list ruleset            # always read back what actually loaded
nft list ruleset | grep -c "ct state"   # must print 0
\end{lstlisting}

\subsection{iptables}

Still widespread, and named explicitly in many site standards. The rules are
a direct transliteration, with one important gap.

\begin{lstlisting}[style=shell,caption={SOURCE host}]
iptables -F && iptables -X
iptables -P INPUT   DROP
iptables -P OUTPUT  DROP
iptables -P FORWARD DROP

iptables -A INPUT  -i lo -j ACCEPT
iptables -A OUTPUT -o lo -j ACCEPT

# Nothing arriving on the diode NIC is ever accepted.
iptables -A INPUT  -i enp1s0 -j DROP

# Exactly one thing may leave by it.
iptables -A OUTPUT -o enp1s0 -p udp -d 192.168.10.2 --dport 6514 -j ACCEPT
iptables -A OUTPUT -o enp1s0 -j DROP

# IPv6 off entirely.
ip6tables -P INPUT DROP; ip6tables -P OUTPUT DROP; ip6tables -P FORWARD DROP
\end{lstlisting}

\begin{lstlisting}[style=shell,caption={DESTINATION host}]
iptables -F && iptables -X
iptables -P INPUT   DROP
iptables -P OUTPUT  DROP
iptables -P FORWARD DROP

iptables -A INPUT  -i lo -j ACCEPT
iptables -A OUTPUT -o lo -j ACCEPT

iptables -A INPUT  -i enp1s0 -p udp -s 192.168.10.1 --dport 6514 -j ACCEPT
iptables -A INPUT  -i enp1s0 -j DROP
iptables -A OUTPUT -o enp1s0 -j DROP

ip6tables -P INPUT DROP; ip6tables -P OUTPUT DROP; ip6tables -P FORWARD DROP
\end{lstlisting}

\begin{warnbox}[iptables cannot filter ARP]
\prog{iptables} operates on IP. ARP traffic passes it untouched, so on
\prog{iptables} alone the \dstside{} host will still answer ARP requests
across the link --- a working, if narrow, return channel. Close it with
\prog{arptables} (or \prog{ebtables}), or use \prog{nftables}, whose
\cmd{netdev} ingress hook covers all of it:
\begin{lstlisting}[style=shell]
arptables -P INPUT DROP
arptables -P OUTPUT DROP
\end{lstlisting}
The \cmd{arp\_ignore=8} sysctl in \S\ref{sec:linklayer} is a second layer for
the same problem, and should be applied regardless.
\end{warnbox}

Note also that \cmd{-P OUTPUT DROP} is host-wide. If the machine has a
management interface, add explicit accepts for it before you apply this over
SSH.

\subsection{pf (OpenBSD and FreeBSD)}

\prog{pf} makes the stateless requirement explicit and readable, and it has a
setting that matters more here than anywhere else: \cmd{set block-policy drop}.
The alternative, \cmd{return}, answers blocked traffic with TCP RSTs and ICMP
unreachables --- packets flowing the wrong way.

\begin{lstlisting}[style=pf,caption={\file{/etc/pf.conf} --- SOURCE host}]
diode_if = "em0"
peer     = "192.168.10.2"
port     = 6514

# Silence. Never answer with an RST or an ICMP unreachable: that is a
# packet travelling from HIGH to LOW.
set block-policy drop
set skip on lo

block all

# The single permitted flow. "no state" is the whole point -- a state
# table entry exists to let the reply back in.
pass out quick on $diode_if inet proto udp \
     from ($diode_if) to $peer port $port no state

block in  quick on $diode_if all
block out quick on $diode_if all
\end{lstlisting}

\begin{lstlisting}[style=pf,caption={\file{/etc/pf.conf} --- DESTINATION host}]
diode_if = "em0"
peer     = "192.168.10.1"
port     = 6514

set block-policy drop
set skip on lo

block all

pass in quick on $diode_if inet proto udp \
     from $peer to ($diode_if) port $port no state

block out quick on $diode_if all
\end{lstlisting}

\begin{lstlisting}[style=shell]
pfctl -nf /etc/pf.conf      # syntax check first
pfctl -f  /etc/pf.conf
pfctl -sr                   # read back the loaded rules
pfctl -si | grep -i state   # state table should stay empty
\end{lstlisting}

\paragraph{OpenBSD versus FreeBSD.} The syntax above works on both.
Differences worth knowing: OpenBSD enables \prog{pf} by default and loads
\file{/etc/pf.conf} at boot, while FreeBSD needs \cmd{pf\_enable="YES"} in
\file{/etc/rc.conf}. FreeBSD's \prog{pf} is a fork of an older OpenBSD
release, so newer OpenBSD-only syntax will not parse there --- everything used
here predates the split. Interface names differ (\cmd{em0}, \cmd{igb0},
\cmd{re0}, \cmd{vtnet0}). FreeBSD users who prefer \prog{ipfw} can express the
same policy with \cmd{ipfw add deny ip from any to any in via em0} on the
\srcside{} host; the stateless requirement is identical --- do not use
\cmd{keep-state} or \cmd{check-state}.

\subsection{A diode in network namespaces, for CI}

The cheapest possible bench: one machine, no hardware, reproducible in a test
script.

\begin{lstlisting}[style=shell]
ip netns add low
ip netns add high
ip link add veth-low type veth peer name veth-high
ip link set veth-low  netns low
ip link set veth-high netns high

ip -n low  addr add 192.168.10.1/24 dev veth-low
ip -n high addr add 192.168.10.2/24 dev veth-high
ip -n low  link set veth-low  up ; ip -n low  link set lo up
ip -n high link set veth-high up ; ip -n high link set lo up

# Static neighbours: no ARP anywhere.
LMAC=$(ip -n low  link show veth-low  | awk '/link\/ether/{print $2}')
HMAC=$(ip -n high link show veth-high | awk '/link\/ether/{print $2}')
ip -n low  neigh replace 192.168.10.2 lladdr "$HMAC" dev veth-low  nud permanent
ip -n high neigh replace 192.168.10.1 lladdr "$LMAC" dev veth-high nud permanent

# Apply the per-side rulesets from above.
ip netns exec low  nft -f /etc/nftables-low.conf
ip netns exec high nft -f /etc/nftables-high.conf

# Run the daemons.
ip netns exec high ./bin/<receiver> ...   # see section 7
ip netns exec low  ./bin/<sender>   ...
\end{lstlisting}

Namespaces add the interface-scoped rules, which is what you want when
testing the firewall policy itself rather than the codec.

\subsection{A unidirectional layer-2 tunnel with TAP}
\label{sec:tap}

Everything so far filters at layer 3 and relies on a ruleset staying in place.
There is a better software construction: bridge a TAP interface to each of two
\emph{real} Ethernet segments and forward frames between the taps with a
program that has no reverse code path. The result is a unidirectional
layer-2 tunnel --- the same Ethernet semantics as the optical diode of
\S\ref{sec:hard}, with software standing in for the missing transmitter.

\begin{figure}[h]
\centering
\begin{tikzpicture}[node distance=7mm]
  %% top row: the SOURCE segment feeds the input tap
  \node[host, minimum width=20mm] (lseg) {LOW\\segment};
  \node[mc, minimum width=14mm, right=of lseg] (e0) {eth0};
  \node[mc, minimum width=18mm, right=of e0, fill=ldblue!8] (brl) {br-low};
  \node[mc, minimum width=19mm, right=of brl, fill=ldamber!12,
        draw=ldamber] (tl) {tapdlow};
  %% bottom row: the output tap feeds the DESTINATION segment (flows right to left)
  \node[host, minimum width=20mm, below=24mm of lseg] (hseg) {HIGH\\segment};
  \node[mc, minimum width=14mm, right=of hseg] (e1) {eth1};
  \node[mc, minimum width=18mm, right=of e1, fill=ldblue!8] (brh) {br-high};
  \node[mc, minimum width=19mm, right=of brh, fill=ldamber!12,
        draw=ldamber] (th) {tapdhigh};
  %% the forwarder, between the two rows
  \node[mc, minimum width=24mm, minimum height=15mm, fill=ldgreen!10,
        draw=ldgreen, right=13mm of $(tl)!0.5!(th)$] (d) {\textbf{tapdiode}};

  \draw[dataflow] (lseg)--(e0); \draw[dataflow] (e0)--(brl);
  \draw[dataflow] (brl)--(tl);
  \draw[dataflow] (tl.east)--(d.north west);
  \draw[dataflow] (d.south west)--(th.east);
  \draw[dataflow] (th)--(brh); \draw[dataflow] (brh)--(e1);
  \draw[dataflow] (e1)--(hseg);

  \node[lbl, below=1.5mm of d, align=center, text=ldgreen!45!black]
       {reads \texttt{in\_fd},\\writes \texttt{out\_fd},\\never the reverse};
  \draw[blocked] (hseg.west) to[out=180,in=180]
        node[midway,fill=white,inner sep=1pt]{\xmark} (lseg.west);
\end{tikzpicture}
\caption{A unidirectional layer-2 tunnel. Each tap is bridged to a real NIC,
         so two physical Ethernet segments are joined --- in one direction.}
\label{fig:tap}
\end{figure}

\subsubsection{Why this is a stronger software construction}

A firewall diode is a \emph{policy} statement: the path exists and a rule
forbids it. This is a \emph{structural} one: the program holds two file
descriptors, reads only the first and writes only the second. The reverse
direction is not disabled or filtered --- it is absent. Nothing to flush,
nothing to misconfigure, and reviewing it means reading eighty lines rather
than auditing a ruleset against a kernel's conntrack behaviour.

\begin{labbox}[It is still software, and still not a security boundary]
The binary can be replaced, \cmd{ptrace}d, or simply bypassed by adding both
NICs to one bridge --- all of which are one command for root on the forwarding
host. It belongs in the same tier as \S\ref{sec:soft}: excellent for a bench,
for CI, and for teaching the layer-2 semantics of a diode. When the guarantee
has to survive a compromise of the host, only optics do that.
\end{labbox}

\subsubsection{The forwarder}

The forwarder is about a hundred lines. It attaches to two existing TAP
devices, then copies frames from the first to the second for as long as it
runs. Table~\ref{tab:tapnotes} covers the choices in it that are not obvious
from the source.

\begin{lstlisting}[style=c,label={lst:tapdiode},caption={%
  \file{tapdiode.c} --- a one-way layer-2 forwarder. Build with
  \cmd{cc -O2 -Wall -Wextra -pedantic -std=c11 -o tapdiode tapdiode.c}}]
/* tapdiode -- a layer-2 one-way bridge between two TAP interfaces.
 *
 * Frames read from the LOW tap are written to the HIGH tap.  The reverse
 * direction is not disabled, commented out, or filtered: it is absent.
 * out_fd is never read from and never appears in a pollfd, so there is no
 * code path along which a frame could travel HIGH -> LOW.
 *
 * Both taps must already exist and be bridged to different Ethernet
 * segments -- see the bridge setup below.
 */
#define _GNU_SOURCE
#include <errno.h>
#include <fcntl.h>
#include <poll.h>
#include <signal.h>
#include <stdio.h>
#include <string.h>
#include <unistd.h>
#include <sys/ioctl.h>

#include <net/if.h>          /* IFNAMSIZ, struct ifreq -- NOT linux/if.h */
#include <linux/if_tun.h>

#define FRAME_MAX 65536      /* jumbo frame + VLAN tags, with room to spare */

static volatile sig_atomic_t stop_requested;

static void on_signal(int sig) { (void) sig; stop_requested = 1; }

/* Attach to the existing TAP device `name'.  Returns a descriptor, or -1. */
static int tap_attach(const char *name)
{
    struct ifreq ifr;
    size_t len = strlen(name);
    int fd;

    /* ifr_name is IFNAMSIZ bytes; validate before copying into it. */
    if (len == 0 || len >= IFNAMSIZ) {
        fprintf(stderr, "tapdiode: bad interface name '%s' (1..%d chars)\n",
                name, IFNAMSIZ - 1);
        return -1;
    }

    fd = open("/dev/net/tun", O_RDWR | O_CLOEXEC);
    if (fd < 0) { perror("tapdiode: open /dev/net/tun"); return -1; }

    memset(&ifr, 0, sizeof ifr);
    /* IFF_TAP   -- layer 2; the buffer is a whole Ethernet frame.
     * IFF_NO_PI -- no 4-byte struct tun_pi prefix to skip. */
    ifr.ifr_flags = IFF_TAP | IFF_NO_PI;
    memcpy(ifr.ifr_name, name, len + 1);

    if (ioctl(fd, TUNSETIFF, &ifr) < 0) {
        fprintf(stderr, "tapdiode: TUNSETIFF %s: %s\n", name, strerror(errno));
        close(fd);
        return -1;
    }
    return fd;
}

int main(int argc, char **argv)
{
    static unsigned char frame[FRAME_MAX];
    unsigned long long forwarded = 0, bytes = 0, lost = 0;
    struct sigaction sa;
    struct pollfd pfd;
    int in_fd, out_fd;

    if (argc != 3) {
        fprintf(stderr, "usage: %s <low-tap> <high-tap>\n"
                        "  Forwards Ethernet frames low -> high only.\n", argv[0]);
        return 2;
    }

    memset(&sa, 0, sizeof sa);
    sa.sa_handler = on_signal;
    if (sigaction(SIGINT, &sa, NULL) < 0 || sigaction(SIGTERM, &sa, NULL) < 0) {
        perror("tapdiode: sigaction");
        return 1;
    }

    if ((in_fd = tap_attach(argv[1])) < 0) return 1;
    if ((out_fd = tap_attach(argv[2])) < 0) { close(in_fd); return 1; }

    fprintf(stderr, "tapdiode: %s -> %s, one way (no reverse code path)\n",
            argv[1], argv[2]);

    /* The input descriptor is the only thing ever polled, and only for
     * readability.  out_fd is write-only by construction. */
    pfd.fd = in_fd;
    pfd.events = POLLIN;

    while (!stop_requested) {
        ssize_t n, w;

        if (poll(&pfd, 1, -1) < 0) {
            if (errno == EINTR) continue;
            perror("tapdiode: poll");
            break;
        }

        n = read(in_fd, frame, sizeof frame);
        if (n < 0) {
            if (errno == EINTR || errno == EAGAIN) continue;
            perror("tapdiode: read");
            break;
        }
        if (n == 0) continue;

        w = write(out_fd, frame, (size_t) n);
        if (w < 0) {
            if (errno == EINTR) continue;
            /* A full queue downstream is loss, which is the honest failure
             * mode of a diode.  Count it; never block, never signal back. */
            lost++;
            continue;
        }
        if (w != n) { lost++; continue; }   /* short write: frame malformed */
        forwarded++;
        bytes += (unsigned long long) n;
    }

    fprintf(stderr, "tapdiode: forwarded %llu frames (%llu bytes), lost %llu\n",
            forwarded, bytes, lost);
    close(in_fd);
    close(out_fd);
    return 0;
}
\end{lstlisting}

\begin{table}[h]
\centering\small
\begin{tabularx}{\textwidth}{@{}p{40mm}X@{}}
\toprule
\textbf{Choice} & \textbf{Reason} \\
\midrule
Only \cmd{in\_fd} is polled &
  \cmd{out\_fd} never appears in a \cmd{pollfd} and \cmd{read()} is never
  called on it, so one-wayness is a property of the control flow --- visible
  to anyone reading the loop, and not dependent on configuration. \\[3pt]
\cmd{poll()} rather than \cmd{select()} &
  No \cmd{FD\_SETSIZE} ceiling, no \cmd{fd\_set} to rebuild each iteration,
  and the intent --- one descriptor, readability only --- is explicit in the
  call. \\[3pt]
\cmd{IFF\_NO\_PI} &
  Suppresses the 4-byte \cmd{struct tun\_pi} prefix, so the buffer holds
  exactly the Ethernet frame and nothing has to be skipped on either
  side. \\[3pt]
Interface names as arguments &
  Bridge membership and addressing can be scripted deterministically instead
  of depending on which \cmd{tapN} the kernel happens to assign. \\[3pt]
Name length checked against \cmd{IFNAMSIZ} &
  \cmd{ifr\_name} is a fixed 16-byte field; the check keeps the copy into it
  bounded. \\[3pt]
\cmd{<net/if.h>}, not \cmd{<linux/if.h>} &
  The glibc and kernel headers define overlapping structures; mixing them
  fails to compile on a current toolchain. \\[3pt]
\cmd{EINTR} retried; every syscall checked &
  A signal must not be mistaken for a closed device, and a failed
  \cmd{read()} must never reach \cmd{write()} as a length. \\[3pt]
Write failure counted, never retried &
  A full queue downstream is loss, and loss is the honest failure mode of a
  diode. Blocking or signalling upstream would both be worse. \\[3pt]
Counters, not a per-frame \cmd{printf} &
  Per-frame output would dominate the runtime; totals are reported once, on
  exit. \\[3pt]
\cmd{O\_CLOEXEC}, signal handling, 64\,KiB buffer &
  Descriptors do not leak across an \cmd{exec}; \cmd{SIGINT}/\cmd{SIGTERM}
  produce a clean summary; jumbo frames and VLAN tags fit with room to
  spare. \\
\bottomrule
\end{tabularx}
\caption{Design notes for Listing~\ref{lst:tapdiode}.}
\label{tab:tapnotes}
\end{table}

\subsubsection{Bridging it to real interfaces}

\begin{lstlisting}[style=shell,caption={On the forwarding host}]
# One tap per side.  Persistent, so they can be bridged before tapdiode runs.
ip tuntap add dev tapdlow  mode tap
ip tuntap add dev tapdhigh mode tap

# One bridge per side.  STP off: BPDUs are bidirectional signalling and have
# no business on a diode.
ip link add br-low  type bridge stp_state 0
ip link add br-high type bridge stp_state 0

ip link set eth0     master br-low     # NIC facing the SOURCE segment
ip link set tapdlow  master br-low
ip link set eth1     master br-high    # NIC facing the DESTINATION segment
ip link set tapdhigh master br-high

for i in eth0 eth1 tapdlow tapdhigh br-low br-high; do ip link set "$i" up; done

# The forwarding host is a pure layer-2 device: give it no address on either
# side, and no IPv6, so it is not itself reachable from either segment.
for i in br-low br-high eth0 eth1 tapdlow tapdhigh; do
    sysctl -qw "net.ipv6.conf.$i.disable_ipv6=1"
done

./tapdiode tapdlow tapdhigh
\end{lstlisting}

The two segments then carry the addressing, and because this is a layer-2
tunnel the hosts on either side are configured exactly as for the optical
diode: same subnet, static neighbour entries, no ARP (\S\ref{sec:linklayer}).

\begin{notebox}[Two bridges, never one]
The whole construction rests on \cmd{br-low} and \cmd{br-high} being separate
bridges. Adding \cmd{eth0} and \cmd{eth1} to the \emph{same} bridge produces
an ordinary bidirectional bridge and silently destroys the property --- and it
is a single command. If you script this, assert the membership afterwards:
\cmd{bridge link show} must never list \cmd{eth0} and \cmd{eth1} under one
master.
\end{notebox}

\begin{notebox}[Frames from HIGH do enter tapdhigh, and stop there]
\cmd{br-high} forwards \dstside{} traffic toward \cmd{tapdhigh}, so those
frames land in that tap's transmit queue --- where nothing ever reads them,
because \cmd{tapdiode} never calls \cmd{read()} on \cmd{out\_fd}. The queue
fills and the kernel discards them. This costs a little memory and no
security: a tap's transmit and receive queues are independent, so an unread
queue cannot exert backpressure on the writes going the other way. If the
waste offends, drop the traffic at the bridge with an \prog{nftables}
\cmd{bridge}-family rule on egress toward \cmd{tapdhigh}.
\end{notebox}

\subsubsection{Verifying it}

Layer 2 needs a layer-2 test; \cmd{ping} proves nothing about frames the IP
stack never generates. Inject a tagged frame with a distinctive EtherType on
one segment and sniff for it on the other:

\begin{lstlisting}[style=shell]
# On the DESTINATION segment: capture (EtherType 0x88b5 is reserved for local use)
tcpdump -ni eth0 -e 'ether proto 0x88b5'

# On the SOURCE segment: inject one frame
python3 - <<'PY'
import socket
s = socket.socket(socket.AF_PACKET, socket.SOCK_RAW); s.bind(('eth0', 0))
s.send(b'\xff'*6 + b'\x02\x00\x00\x00\x00\x01' + b'\x88\xb5' + b'DIODETEST')
PY
\end{lstlisting}

The frame must appear on the \dstside{} segment. Then reverse the two
commands: injecting on \dstside{} must produce \emph{nothing} on \srcside{},
while a capture on \cmd{br-high} confirms the frame really was injected --- so
that a silent \srcside{} means one-wayness and not a broken injector.

\begin{tipbox}
Sniff with \cmd{socket.htons(3)} if you write the capture side in Python. The
\cmd{AF\_PACKET} protocol argument is in network byte order and Python does
not convert it, so a literal \cmd{3} binds to EtherType \cmd{0x0300} and sees
nothing --- a mistake that looks exactly like a failed diode.
\end{tipbox}

%% =====================================================================
\section{Hardware diodes}
\label{sec:hard}

\subsection{What a media converter is, and which ones work}

A media converter bridges copper Ethernet and optical fiber. Internally it is
two transceivers back to back: copper in $\rightarrow$ fiber out, fiber in
$\rightarrow$ copper out. The optical side is what interests us, and only one
kind of optical side can be used.

\begin{table}[h]
\centering\small
\begin{tabularx}{\textwidth}{@{}p{40mm}p{20mm}X@{}}
\toprule
\textbf{Optical interface} & \textbf{Usable?} & \textbf{Why} \\
\midrule
Dual-fiber, discrete TX and RX ports (two SC or two LC ferrules) &
  \textbf{Yes} &
  TX and RX are physically separate. You can connect one and not the other ---
  which is the entire mechanism of an optical diode. \\[3pt]
Single-fiber bidirectional (BiDi / WDM, one strand, e.g.\ \nm{1310} out and
\nm{1550} back) &
  \textbf{No} &
  Both directions share one connector and one strand, separated only by
  wavelength inside the module. There is no TX-only path to expose. \\
\bottomrule
\end{tabularx}
\caption{Only two-strand converters with separate TX and RX ports can build a
         diode. Single-port BiDi models cannot, at any price.}
\end{table}

The obstacle is that a converter with nothing on its RX port sees no carrier,
declares the fiber link down, and stops forwarding. The paper records this
directly: converters with a disconnected RX port ``will automatically turn off
unless electronic hacks are in place'' \cite{anton2024}. So a diode is not
built by omitting a cable. It is built by feeding the transmitter light that
carries no information from the far side. There are exactly two ways to do
that, and they are the next two subsections.

\subsection{Topology A --- two converters and a fiber splitter}
\label{sec:topoA}

Split the transmitter's own outgoing light and feed half of it back into the
transmitter's own receiver. The transmitter sees carrier --- its own --- and
holds the link up. The other half of the light is the data path.

\begin{figure}[h]
\centering
\resizebox{\textwidth}{!}{%
\begin{tikzpicture}
  \node[host, minimum width=20mm] (low) {SOURCE host};
  \node[mc, minimum width=20mm, right=13mm of low]
       (A) {\textbf{MC-A}\\\scriptsize TX \quad RX};
  \node[splitter, right=13mm of A] (S) {1$\times$2\\splitter};
  \node[mc, minimum width=20mm, right=13mm of S]
       (B) {\textbf{MC-B}\\\scriptsize TX \quad RX};
  \node[host, minimum width=20mm, right=13mm of B] (high) {DESTINATION host};

  \draw[copper] (low) -- (A);
  \draw[copper] (B) -- (high);
  \draw[fiber] (A.east) -- node[lblg,above]{TX} (S.west);
  \draw[fiber] (S.east) -- node[lblg,above]{$-$\dbb{4}} node[lblg,below]{data} (B.west);
  \draw[fiber] (S.south) to[out=-90,in=-90]
        node[lblg,below]{$-$\dbb{4}\ \ return carrier (its own echo)} (A.south);

  \node[lblr, above=6mm of B, align=center, anchor=south]
       {MC-B fiber TX:\\\textbf{leave unconnected, capped}};
  \draw[ldred,dashed] (B.north) -- +(0,5mm);
\end{tikzpicture}}
\caption{Two-converter diode. The splitter is passive and reciprocal, so the
         unused legs matter as much as the used ones.}
\label{fig:splitter}
\end{figure}

\begin{warnbox}[A splitter is reciprocal]
An optical splitter is a passive, bidirectional device. Light entering any leg
emerges from the others. In Figure~\ref{fig:splitter} the only reason no
information flows backwards is that \textbf{MC-B's transmitter is not
connected to anything}. Cap that port. If someone later patches MC-B's TX into
a spare splitter leg --- or into the return fiber --- \dstside{} data reaches
MC-A's receiver and the diode is gone, with no visible change to the front
panel LEDs.
\end{warnbox}

\begin{notebox}[The transmitter hears its own echo]
MC-A receives its own transmission, so MC-A's copper port re-emits every frame
the \srcside{} host sent. This is harmless to the diode property --- nothing
from \dstside{} is involved --- but it has two practical consequences:
\begin{itemize}
\item \textbf{Never put a switch between the \srcside{} host and MC-A.} The
      echo becomes a forwarding loop and then a broadcast storm. Connect the
      host's NIC directly.
\item Disable ARP and IPv6 on the diode interface (\S\ref{sec:linklayer}) so
      the host does not react to its own reflected broadcasts.
\end{itemize}
Topology B does not have this behaviour, which is a quiet point in its favour.
\end{notebox}

\subsection{Sourcing the splitter, and the power budget}
\label{sec:budget}

The part itself is unremarkable: an ordinary passive $1\times2$ 50:50
splitter. It works with all three converters below, including the \nm{850}
multimode MC200CM.\footnote{The \emph{Sensors} paper states that this approach
``did not work'' for the MC200CM, and reports the three-converter arrangement
being used instead (\cite{anton2024}, \S2.2). That wording describes the parts
available at the time rather than a limitation of the method: no \nm{850}
multimode splitter could be obtained within the timeframe of that work. The
one eventually used had to be made to special order and took months to
arrive, and with it in hand the two-converter topology works on the MC200CM
exactly as it does at \nm{1310}.} Topology A is not limited by physics on any
of them.

What differs between wavelengths is \textbf{availability and price}, and the
gap is large. Splitters at \nm{1310} are a fiber-to-the-home product ---
single-mode, specified across \nm{1260}--\nm{1650}, manufactured by the
million, stocked everywhere and cheap. Multimode splitters at \nm{850} are a
specialty item: the \nm{850} splitter used for the gigabit build here had to
be made to special order and took months to arrive. The component works
perfectly well once you have it; the difficulty is entirely in getting one.

\begin{tipbox}[What this means for a build]
Choose Topology A freely at \nm{1310} (MC100CM or MC210CS) --- the splitter is
a catalogue part you can have tomorrow. At \nm{850} (MC200CM) Topology A works
just as well, but budget for a special order and a long lead time. If you need
the diode sooner than the splitter can arrive, Topology B
(\S\ref{sec:topoB}) needs no splitter at all and uses a third converter you
can buy off the shelf. Match the splitter's fiber type to the link:
single-mode splitter on single-mode fiber, multimode on multimode.
\end{tipbox}

\subsubsection{Power budget and attenuation}

A 50:50 split costs $10\log_{10}2 = 3.01$\,dB per leg by construction, plus a
little excess loss and connector loss --- call it \dbb{4} in round numbers.
Every class below has room for that, which is why the splitter build works
across the range. The figures are worth having anyway, because they tell you
when to \emph{reduce} the light rather than worry about losing it:

\begin{table}[h]
\centering\small
\begin{tabularx}{\textwidth}{@{}llcccccX@{}}
\toprule
\textbf{Model} & \textbf{Class} & \textbf{$\lambda$} &
\textbf{TX min} & \textbf{TX max} & \textbf{RX sens.} & \textbf{RX max} &
\textbf{Budget} \\
\midrule
MC100CM & 100BASE-FX  & \nm{1310} & \dbm{-20}   & \dbm{-14} & \dbm{-29} &
  \dbm{-14} & $\approx$\dbb{9} \\
MC200CM & 1000BASE-SX & \nm{850}  & \dbm{-9.5}  & \dbm{0}   & \dbm{-17} &
  \dbm{0}   & $\approx$\dbb{7.5} \\
MC210CS & 1000BASE-LX & \nm{1310} & \dbm{-11}   & \dbm{-3}  & \dbm{-19} &
  \dbm{-3}  & $\approx$\dbb{8} \\
\bottomrule
\end{tabularx}
\caption{Figures for the \emph{optical class}, not for these units:
         \textbf{TP-Link publishes no optical power data for the MC100CM,
         MC200CM or MC210CS} --- their datasheets give only wavelength, fiber
         type and distance. These are the published class specifications
         \cite{ieee8023}; a particular module often does better. ``Budget''
         is worst-case launch power minus worst-case receiver sensitivity.
         All three have room for a $\approx$\dbb{4} splitter.}
\end{table}

Note the pattern in the two maximum columns: in every class the maximum launch
power and the maximum receive power are the same number
(\dbm{-14}, \dbm{0}, \dbm{-3}). That is deliberate, and it is what the
attenuation advice below turns on.

\begin{tipbox}[When to attenuate]
A receiver has a maximum input power as well as a sensitivity floor, and
exceeding it saturates the front end and raises the bit error rate --- which
reaches the application as CRC failures and dropped frames.

Within one optical class this is designed out: maximum launch power equals
maximum receive power in every class in the table above (\dbm{-14} for
100BASE-FX, \dbm{0} for 1000BASE-SX, \dbm{-3} for 1000BASE-LX), so a
compliant transmitter cannot overload a compliant receiver of the same class
even over a metre of fiber. Three matched converters on a bench need no
attenuation.

The case that does need it is the paper's: a \emph{long-range} converter or
SFP --- one built to drive tens of kilometres, and launching accordingly ---
used across the short fiber a diode actually needs \cite{anton2024}. Fixed
attenuators (\dbb{3}, \dbb{5}, \dbb{10}) are inexpensive; size them from the
launch power and the receiver's maximum in your datasheets rather than by
trial. They are also how you trim one leg of a splitter without touching the
other.
\end{tipbox}

\subsection{Topology B --- three converters in a triangle}
\label{sec:topoB}

No splitter to source, and no echo. A third converter acts purely as
a carrier source. Its copper port is left unplugged, so it emits an idle
optical carrier and can never carry data.

\begin{figure}[h]
\centering
\resizebox{\textwidth}{!}{%
\begin{tikzpicture}
  \node[host, minimum width=20mm] (low) {SOURCE host};
  \node[mc, minimum width=20mm, right=15mm of low] (A) {\textbf{MC-A}\\\scriptsize input};
  \node[mc, minimum width=20mm, right=26mm of A] (B) {\textbf{MC-B}\\\scriptsize output};
  \node[host, minimum width=20mm, right=15mm of B] (high) {DESTINATION host};
  \node[mc, minimum width=34mm, below=20mm of $(A)!0.5!(B)$,
        fill=ldamber!10, draw=ldamber]
       (C) {\textbf{MC-C} carrier source\\\scriptsize copper port UNPLUGGED};

  \draw[copper] (low) -- (A);
  \draw[copper] (B) -- (high);
  \draw[fiber] (A.north east) to[out=45,in=135]
        node[lblg,above]{\textbf{1.} MC-A TX $\rightarrow$ MC-B RX \ \emph{(data)}} (B.north west);
  \draw[fiber] (B.south) to[out=-90,in=0]
        node[lblg,right,pos=0.3]{\textbf{2.} MC-B TX $\rightarrow$ MC-C RX} (C.east);
  \draw[fiber] (C.west) to[out=180,in=-90]
        node[lblg,left,pos=0.7]{\textbf{3.} MC-C TX $\rightarrow$ MC-A RX} (A.south);
  \node[lblr, below=1.5mm of C, align=center]
       {copper unplugged $\Rightarrow$ links 2 and 3 carry carrier, never data};
\end{tikzpicture}}
\caption{Three-converter diode. The security argument rests on one unplugged
         copper port.}
\end{figure}

\paragraph{Why it is safe.} Trace a hypothetical leak. \dstside{} data enters
MC-B's copper port, leaves MC-B's fiber TX on link~2, and arrives at MC-C's
fiber RX. MC-C converts it to copper --- and its copper port goes nowhere, so
the frames are discarded. MC-C's fiber TX is driven by MC-C's \emph{copper
receiver}, which has no link and no traffic, so link~3 carries idle carrier
only. The chain is broken at MC-C's unplugged copper port, and nowhere else.

\paragraph{Why link 2 exists at all.} MC-C needs carrier on its own RX before
it will bring its optical side up and start transmitting. Link~2 supplies it.
The paper states the same requirement: the connection from the output
converter to the signal emulator ``is necessary in order to keep the signal
emulator running'' \cite{anton2024}.

\begin{warnbox}[The single point of failure is a copper port]
Everything rests on MC-C's copper port being unplugged. Cap it, label it
\emph{``MUST REMAIN DISCONNECTED --- DATA DIODE''}, and put it first on the
audit checklist in \S\ref{sec:commission}. A well-meaning administrator
patching that port into the \srcside{} switch converts the diode into an
ordinary bidirectional link, and no LED changes colour to say so.
\end{warnbox}

\begin{tipbox}[If MC-C will not light up]
Some converters implement link-fault pass-through and go dark on fiber when
their copper link is down. If MC-C behaves that way, give it a copper link
that leads nowhere: an isolated switch port on a dead VLAN with no uplink, or
a loopback plug. Never a port on the \srcside{} network --- that would rebuild
the return path you just broke.
\end{tipbox}

\subsection{Choosing between them}

Topology B is the default. It costs one more converter and one more power
Both work at every wavelength here. Topology A is the tidier build --- one
box fewer, one power supply fewer, and a passive component in place of an
active one. Topology B trades that for parts you can buy off the shelf
anywhere, and it avoids the transmitter echo. At \nm{1310} the choice is
genuinely free; at \nm{850} it comes down to whether you can wait for the
splitter.

\begin{table}[h]
\centering\small
\begin{tabularx}{\textwidth}{@{}p{38mm}XX@{}}
\toprule
 & \textbf{A: 2 converters + splitter} & \textbf{B: 3 converters} \\
\midrule
Parts            & 2 converters, 1 splitter, 3 patch cables &
                   3 converters, 3 patch cables \\
Power supplies   & 2 & 3 \\
Optical margin   & Reduced by $\approx$\dbb{4}; within budget on all three
                   models &
                   Full budget on every link \\
Procurement      & Trivial at \nm{1310}; the \nm{850} splitter is a special
                   order with a long lead time &
                   Everything is catalogue stock at any wavelength \\
Transmitter echo & Yes --- host hears itself; no switch on the \srcside{} leg &
                   No \\
Failure points   & 2 active + 1 passive & 3 active \\
Critical rule    & MC-B's TX must stay unconnected &
                   MC-C's copper must stay unplugged \\
\midrule
\textbf{Use when} & You have the splitter, and want fewer boxes &
                   You want to build today, at any wavelength \\
\bottomrule
\end{tabularx}
\end{table}

\subsection{Power, environment and maintenance}
\label{sec:maint}

Drawn from the paper's \S4.3, which is the most practically useful part of it
\cite{anton2024}:

\begin{itemize}
\item \textbf{Power.} These converters run on 9\,V DC via a
      centre-positive barrel connector; the MC200CM draws about 600\,mA. The
      supplies shipped with them are unremarkable. A single stabilised
      laboratory supply --- the paper uses an XP Power ECE60US09-S, 9\,V at
      6.67\,A --- will run six diodes at once and is far better regulated,
      though it costs an order of magnitude more than a converter.
\item \textbf{UPS.} A diode is on the critical path for log delivery and
      draws very little. Put it on a UPS.
\item \textbf{Cooling.} Passive. Six diodes fit in a 2U case with no fans.
\item \textbf{Attenuators} for short runs with long-range optics
      (\S\ref{sec:budget}).
\item \textbf{MTBF} is comparable to a network card, but a three-converter
      diode has three of them plus their supplies. Keep a spare set.
\item \textbf{Latency} is negligible --- nanoseconds. It never appears in
      end-to-end measurements.
\item \textbf{Managed converters} offer SNMP, fault alerts and test modes, and
      are a genuine security risk on a diode: a management plane is a second
      network path into the boundary device. The paper deliberately uses
      unmanaged units. If you use managed ones, be certain their management
      interfaces are on an isolated VLAN, and understand that you have added
      software to a control whose entire value was that it had none.
\end{itemize}

\subsection{Scaling}

Diodes are cheap, so scale out rather than up. Run several in parallel
between the same pair of networks, using separate NICs or IP aliases on both
sides, and spread messages across them round-robin --- the paper's
Figure~10(a) arrangement. The receiver binds the same port on every address,
so the links aggregate transparently. This is also the answer when the paper's
repeat factor would otherwise need to exceed about 50: add a second diode
rather than more redundancy on one.

\subsection{Embedded scale: a three-wire SPI diode}
\label{sec:spi}

Media converters are the right answer between two networks. Between two
microcontrollers --- a sensor gateway that must report into a protected
monitoring domain, a safety PLC that must emit telemetry but never accept a
command --- they are absurd: three converters, three power supplies and fiber
patch cords to join two boards that sit ten centimetres apart.

At that scale the diode is three wires. SPI already keeps its two data
directions on \emph{physically separate conductors}, so removing one of them
removes a direction. This is the same argument as the optical diode, at chip
scale and for a few cents.

\begin{figure}[h]
\centering
\begin{tikzpicture}[node distance=34mm]
  \node[host, minimum width=32mm, minimum height=22mm]
       (lo) {\textbf{LOW} MCU\\\scriptsize SPI \textbf{master}};
  \node[host, minimum width=32mm, minimum height=22mm, right=of lo]
       (hi) {\textbf{HIGH} MCU\\\scriptsize SPI \textbf{slave}};

  \draw[dataflow] ([yshift=6mm]lo.east)  -- node[lbl,above]{SCLK}
                  ([yshift=6mm]hi.west);
  \draw[dataflow] ([yshift=0mm]lo.east)  -- node[lbl,above]{MOSI}
                  ([yshift=0mm]hi.west);
  \draw[copper]   ([yshift=-6mm]lo.east) -- node[lbl,above]{GND}
                  ([yshift=-6mm]hi.west);

  \draw[blocked] ([yshift=-16mm]hi.west) -- ([yshift=-16mm]lo.east)
                 node[midway,fill=white,inner sep=1.5pt]{\xmark}
                 node[lblr,midway,below=3.5mm]{MISO --- conductor not fitted};
\end{tikzpicture}
\caption{A three-wire SPI diode. Every conductor is driven by the \srcside{}
         side and read by the \dstside{} side; the return data line is not
         fitted.}
\label{fig:spi}
\end{figure}

\subsubsection{The sender must be the master}

This is the design decision that matters, and it is not the obvious one.

An SPI master drives SCLK, CS and MOSI, and reads MISO. A slave drives MISO
and reads the rest. For information to flow \srcside{}~$\rightarrow$~\dstside{}
with \emph{no} conductor driven from \dstside{}, the \srcside{} side must be
the master: it drives clock and data, and the \dstside{} side only ever
listens.

The tempting alternative --- make the \dstside{} side the master so it
``pulls'' data --- quietly reintroduces a channel. A master drives the clock,
so SCLK would run \dstside{}~$\rightarrow$~\srcside{}, and the \dstside{}
side could then modulate the timing or gating of clock edges into a signal the
\srcside{} side can observe. It is a low-rate channel, but it is a real one,
and it exists by construction rather than by bug.

\begin{table}[h]
\centering\small
\begin{tabularx}{\textwidth}{@{}llX@{}}
\toprule
\textbf{Signal} & \textbf{Fitted?} & \textbf{Direction and role} \\
\midrule
SCLK & yes & \srcside{}~$\rightarrow$~\dstside{}. Clock; the sender sets the
             pace. \\
MOSI & yes & \srcside{}~$\rightarrow$~\dstside{}. The data. \\
GND  & yes & Common return. See the isolation note below. \\
CS   & optional & \srcside{}~$\rightarrow$~\dstside{}. A fourth conductor
             that gives hardware frame boundaries; strongly recommended. \\
MISO & \textbf{no} & Would be \dstside{}~$\rightarrow$~\srcside{}. Omit the
             conductor entirely --- do not merely leave it unconfigured. \\
\bottomrule
\end{tabularx}
\caption{Three conductors, or four with CS. The diode property is the row
         that is not there.}
\end{table}

\begin{warnbox}[``3-wire SPI mode'' in a datasheet usually means the opposite]
Many MCU and sensor datasheets use ``3-wire SPI'' for \emph{half-duplex} SPI,
in which MOSI and MISO are collapsed onto one \textbf{bidirectional} data pin
(often \cmd{SDIO} or \cmd{SDA}). That is a shared, driven-from-both-ends
conductor --- precisely what a diode must not have. What is wanted here is
ordinary full-duplex SPI with the MISO wire physically absent. Read the pin
direction, not the marketing name.
\end{warnbox}

\begin{warnbox}[Never add a ready line]
Without MISO there is no flow control: the master clocks bits out whether or
not the slave is keeping up. The standard embedded fix is a GPIO from slave to
master meaning ``ready'' or ``busy'' --- and that is a wire from \dstside{} to
\srcside{}. Adding it destroys the diode as thoroughly as fitting MISO would,
and it is the single most likely mistake in this design because every SPI
tutorial recommends it. Handle overrun the way the rest of this handbook
handles loss: clock slowly enough that the receiver always keeps up, use DMA
with double buffering, and send redundancy in advance.
\end{warnbox}

\subsubsection{Isolation}

A shared ground is still a conductor between the two domains. Where that
matters --- different power domains, an untrusted \srcside{} board, or simply a
wish to make the property auditable at a glance --- put a digital isolator or
an optocoupler on each of the three signals. This is worth doing for its own
sake:

\begin{itemize}
\item An optocoupler is an LED illuminating a phototransistor. It is
      \emph{physically} incapable of conducting backwards, so each isolated
      line enforces its own direction in hardware --- the optical diode
      argument, in a DIP-4 package.
\item Galvanic isolation removes the shared ground return, so the two domains
      no longer share a reference at all.
\item Check the part's data rate. Ordinary optocouplers manage a few hundred
      kbit/s; use a proper digital isolator for anything faster, and confirm
      it is a single-channel unidirectional device rather than a
      bidirectional or mixed-direction part.
\end{itemize}

\subsubsection{Framing, and what to run over it}

With CS fitted, the slave gets frame boundaries in hardware and resynchronizes
for free. Without it the receiver must find its own boundaries in the bit
stream, so frame in band: a sync word, a length, the payload, and a check
value. That is the same bounded-cursor discipline as everywhere else in this
system --- a length is validated against what remains before it is ever used as
an index, and a length that lies ends the frame instead of becoming a pointer.

Everything the rest of this handbook says about a feedback-free link applies
unchanged: no acknowledgement, so no retransmission; redundancy must be sent
ahead of time; and authentication is what distinguishes a corrupted frame from
a forged one.

\begin{notebox}[Reusing the proven core on an MCU]
The SPARK units in \file{src/core} are fixed-capacity and allocation-free by construction, so
they do not depend on an operating system and can be cross-compiled for a bare
metal target. \amp{} and \deamp{} themselves cannot: they are Linux daemons
built on UDP sockets, and the sockets would have to be replaced by SPI
transfers.

The capacities are the obstacle worth planning for. As analyzed, the instance
assumes \cmd{Secure.Max\_Plain}~=~65536, \cmd{Syslog.Max\_Record}~=~65535 and a
relay holding sixteen in-flight messages --- megabytes of static buffers, which
no small microcontroller has. Reducing them is a supported change, but it
changes what was proved: the capacities are part of the analyzed instance, so
a shrunken build must be re-proved before any of the assurance claims in
\S\ref{sec:threat} carry over to it.
\end{notebox}

\subsubsection{Commissioning an SPI diode}

\begin{enumerate}
\item \textbf{Count the wires.} Three, or four with CS. This is the whole
      audit, and it is why the design is attractive --- correctness is visible
      without instruments.
\item \textbf{Continuity.} With both boards powered down, confirm an open
      circuit between the \dstside{} MCU's MISO pin and every conductor in
      the cable. Repeat after any rework.
\item \textbf{Drive test.} Configure the \dstside{} MCU to drive its MISO pin
      hard, high and low, in a loop. The \srcside{} MCU must see nothing on
      any pin. This catches a MISO track left on a shared PCB, which a wire
      count will not.
\item \textbf{Isolator orientation.} If optocouplers are fitted, confirm each
      one's input side faces \srcside{}. A part installed backwards will
      usually appear to work --- in the wrong direction.
\item \textbf{Overrun.} Run the sender at full rate for an extended period and
      confirm the receiver's dropped-frame counter behaves as loss, not as
      corruption reaching the application.
\end{enumerate}

%% =====================================================================
\section{Commissioning and verification}
\label{sec:commission}

A diode you have not tested is a diode you have assumed. The tests below are
ordered so that each one fails cheaply before the next.

\subsection{Physical audit}

\begin{enumerate}
\item Trace every fiber by hand against the wiring table for your topology.
\item Topology B: confirm MC-C's copper port is unplugged and capped.
      Topology A: confirm MC-B's fiber TX is unconnected and capped.
\item Confirm the \srcside{} host connects \emph{directly} to MC-A, with no
      switch (mandatory for Topology A because of the echo).
\item Photograph the finished wiring. Attach it to the commissioning record;
      it is what a future auditor compares against.
\end{enumerate}

\subsection{Carrier independence}

Unplug MC-B's copper cable from the \dstside{} host. MC-A must keep its fiber
link LED. If MC-A drops, its carrier is arriving from the far end and the link
is bidirectional.

Repeat in the other direction: power MC-C down (Topology B). MC-A should lose
its fiber link, confirming that MC-C --- and not MC-B --- is its carrier
source.

\subsection{Active reverse probing}

Capture everything on the \srcside{} host for the duration:

\begin{lstlisting}[style=shell]
# SOURCE host -- runs for the whole test, must stay silent
tcpdump -ni enp1s0 -e -vv -w /tmp/diode-commissioning.pcap
\end{lstlisting}

From the \dstside{} host, attempt every reasonable back channel:

\begin{lstlisting}[style=shell]
ping   -c 5            192.168.10.1        # ICMP echo
arping -c 5 -I enp1s0  192.168.10.1        # link layer
nc -u -w2              192.168.10.1 9999   # UDP
nc    -w5              192.168.10.1 22     # TCP -- must time out
hping3 -S -c 5 -p 22   192.168.10.1        # raw SYN
hping3 -1 -c 5         192.168.10.1        # raw ICMP
\end{lstlisting}

Then verify the capture contains nothing whose source is the \dstside{} host:

\begin{lstlisting}[style=shell]
tcpdump -nr /tmp/diode-commissioning.pcap "src host 192.168.10.2"
# must print no packets
\end{lstlisting}

\subsection{Forward path and end-to-end}

\begin{lstlisting}[style=shell]
# DESTINATION host
tcpdump -ni enp1s0 udp port 6514
# SOURCE host
echo hello | nc -u -w1 192.168.10.2 6514

# then, with the daemons running:
# SOURCE host
logger -n 127.0.0.1 -P 1514 -p local0.crit "diode commissioning $(date -Is)"
# DESTINATION host
journalctl -f | grep "diode commissioning"
\end{lstlisting}

Finally, pull MC-A's TX fiber: the \dstside{} side must stop receiving
immediately, and the \srcside{} side must be entirely unaffected --- it has no
way to notice.

\subsection{Recurring audit}

There is no remote diagnostic and no alarm. Failure and tampering are both
silent, so both need a calendar entry.

\begin{table}[h]
\centering\small
\begin{tabularx}{\textwidth}{@{}p{26mm}X@{}}
\toprule
\textbf{Interval} & \textbf{Check} \\
\midrule
Continuous & \dstside{}-side monitoring alerts if the log stream stops. This
             is the only automatic failure indication you get --- configure it. \\
Monthly    & Physical inspection: all LEDs as expected, MC-C's copper port
             still unplugged, seals intact, photo still matches. \\
Quarterly  & Repeat the active reverse probe of \S6.3 in full. \\
On change  & Any rewiring, converter replacement, or firmware change re-runs
             the entire commissioning sequence from \S6.1. \\
\bottomrule
\end{tabularx}
\end{table}

%% =====================================================================
\section{Running opc-diode across the link}
\label{sec:integrate}

The relay is two programs. Start the receiver first, so nothing is lost.

\begin{lstlisting}[style=shell,caption={DESTINATION side --- recover and re-emit to subscribers}]
./bin/od_receiver <diode_port> <out_ip> <out_port> [--key HEX64]

# e.g. recover on 9702, re-emit NetworkMessages to subscribers on 127.0.0.1:9703
./bin/od_receiver 9702 127.0.0.1 9703 --key <64 hex>
\end{lstlisting}

\begin{lstlisting}[style=shell,caption={SOURCE side --- take the publisher's output, protect, transmit}]
./bin/od_sender <in_port> <diode_ip> <diode_port> [--parity M] [--interleave D]
                [--pace-us N] [--key HEX64]

# e.g. take NetworkMessages on 9701, add 3 parity fragments, blast to the diode
./bin/od_sender 9701 <receiver_ip> 9702 --parity 3 --pace-us 200 --key <64 hex>
\end{lstlisting}

The publisher points its PubSub output at the sender's input port; subscribers
listen on the receiver's output port. One-way throughout.

\subsection{The options that matter}

\begin{description}[leftmargin=0pt, style=nextline]
\item[\cmd{--parity M} --- how much loss you survive.]
  $M$ parity fragments are added to the $K$ data fragments of a message, and
  any $K$ of the $K+M$ rebuild it. Raise it first when messages go missing.
\item[\cmd{--interleave D} --- what \emph{shape} of loss you survive.]
  Spreads the fragments of $D$ messages in time (packet 1 of each, then packet
  2 of each), so a burst becomes a recoverable one-per-message trickle rather
  than the loss of one whole message. Costs latency, not bandwidth.
\item[\cmd{--pace-us N} --- how hard you drive the link.]
  Microseconds between packets. A diode has no flow control, so pacing is what
  stops the sender overrunning the receiver.
\item[\cmd{--key HEX64} --- authentication.]
  ChaCha20-Poly1305, on \emph{both} ends. The sender seals each NetworkMessage
  before fragmenting; the receiver verifies the tag and drops anything that
  fails. Without it, payloads cross in the clear. The key is pre-shared out of
  band --- see the note in \S\ref{sec:threat} on why no key agreement is
  possible here.
\end{description}

opc-diode also offers a DPDK kernel-bypass transport (\cmd{--with-dpdk} and
\cmd{--eal}); see the project README, which is the authority on the CLI. The
assurance argument, including what is proved and what is trusted, is in
\file{docs/ASSURANCE.md}.

%% =====================================================================
\section{Troubleshooting}

% longtable has no X column type; widths are explicit (text block is 160mm).
\begin{longtable}{@{}p{40mm}p{34mm}p{73mm}@{}}
\toprule
\textbf{Symptom} & \textbf{Likely cause} & \textbf{Check} \\
\midrule
\endhead
No fiber link LED on MC-A &
  No carrier on its RX &
  Topology B: is MC-C powered and lit? Is link~3 seated? Topology A: is the
  return leg of the splitter connected? \\
MC-C completely dark (Topology B) &
  Link-fault pass-through, copper down &
  Give MC-C a dead-VLAN copper link or loopback plug (\S\ref{sec:topoB}). \\
Link comes up, then flaps &
  Contamination, a bad splice, or a marginal connector &
  Clean and reseat the ferrules first --- this is the usual cause by a wide
  margin. Only suspect overload if you are running long-range optics over a
  short fiber (\S\ref{sec:budget}). \\
\dstside{} receives nothing, LEDs all good &
  ARP or firewall &
  \cmd{ip neigh show} must report PERMANENT. Check \cmd{rp\_filter} if the
  two sides are on different subnets. \cmd{tcpdump} on the \dstside{} NIC. \\
\dstside{} sees datagrams but produces no output &
  Tag verification failing &
  Keys differ between the two sides, or one side has \cmd{--key} and the
  other does not. \\
High loss under load &
  Receiver overrun &
  Raise \cmd{--parity}, add \cmd{--pace-us}, raise the socket receive buffer,
  or add a second diode in parallel. \\
Loss in bursts, fine otherwise &
  No interleaving &
  Raise \cmd{--interleave}; this is the case it exists for. \\
Whole groups of lines vanish together &
  Batch coupling &
  A failed block loses all $N$ records. Lower \cmd{--batch} or raise
  \cmd{--parity}. \\
Broadcast storm on the \srcside{} segment &
  Topology A echo through a switch &
  Connect the \srcside{} host directly to MC-A (\S\ref{sec:topoA}). \\
\srcside{} host logs traffic from itself &
  Topology A echo, expected &
  Harmless. Confirm the source is the \srcside{} host, not \dstside{} --- if
  it is \dstside{}, stop and re-run \S\ref{sec:commission}. \\
\bottomrule
\end{longtable}

%% =====================================================================
\begin{thebibliography}{99}
\small
\bibitem{anton2024}
A.-A. Anton, P. Csereoka, E. A. Capot\u{a} and R.-D. Cioarg\u{a},
``Enhancing Syslog Message Security and Reliability over Unidirectional Fiber
Optics'', \emph{Sensors}, vol.~24, no.~20, 6537, 2024.
\url{https://doi.org/10.3390/s24206537}

\bibitem{bell1976}
D. E. Bell and L. J. LaPadula,
``Secure Computer System: Unified Exposition and Multics Interpretation'',
MITRE Technical Report ESD-TR-75-306, 1976.

\bibitem{rfc5424}
R. Gerhards, ``The Syslog Protocol'', RFC 5424, IETF, 2009.

\bibitem{rfc8439}
Y. Nir and A. Langley, ``ChaCha20 and Poly1305 for IETF Protocols'',
RFC 8439, IETF, 2018.

\bibitem{ieee8023}
IEEE Std 802.3, \emph{IEEE Standard for Ethernet}: clause 26 and the FDDI PMD
it adopts (100BASE-FX), and clause 38 (1000BASE-SX, 1000BASE-LX). Source of
the per-class optical figures in \S\ref{sec:budget}. Those figures are class
specifications: the vendor datasheets for the MC100CM, MC200CM and MC210CS
give wavelength, fiber type and distance only, and publish no optical power
data at all.

\bibitem{etherled}
M. Guri, ``ETHERLED: Sending Covert Morse Signals from Air-Gapped Devices via
Network Card (NIC) LEDs'', \emph{IEEE International Conference on Cyber
Security and Resilience (CSR)}, 2022, pp.~163--170.

\bibitem{lantenna}
M. Guri, ``LANTENNA: Exfiltrating Data from Air-Gapped Networks via Ethernet
Cables Emission'', \emph{IEEE 45th Annual Computers, Software, and
Applications Conference (COMPSAC)}, Madrid, 2021, pp.~745--754.

\bibitem{powerhammer}
M. Guri, B. Zadov, D. Bykhovsky and Y. Elovici, ``PowerHammer: Exfiltrating
Data From Air-Gapped Computers Through Power Lines'', \emph{IEEE Transactions
on Information Forensics and Security}, vol.~15, pp.~1879--1890, 2020.

\bibitem{gsmem}
M. Guri, A. Kachlon, O. Hasson, G. Kedma, Y. Mirsky and Y. Elovici,
``GSMem: Data Exfiltration from Air-Gapped Computers over GSM Frequencies'',
\emph{24th USENIX Security Symposium}, Washington DC, 2015, pp.~849--864.

\bibitem{sparknacl}
R. Chapman, \emph{SPARKNaCl}: a verified Ada/SPARK implementation of TweetNaCl.
\url{https://github.com/rod-chapman/SPARKNaCl}

\bibitem{assurance}
opc-diode project, \emph{Trusted-boundary assurance argument},
\file{docs/ASSURANCE.md}.
\end{thebibliography}

\vfill
\begin{center}\small
\textcolor{ldgray}{%
Made \copyleft{} libre (free as in freedom) at
Politehnica University of Timi\c{s}oara.\\
Copyright \copyright{} 2026 Alin-Adrian Anton.\\
This document is licensed under the Creative Commons Attribution 4.0
International Licence (CC~BY~4.0) or any later version.\\
\url{https://creativecommons.org/licenses/by/4.0/}}
\end{center}

\end{document}
